\documentclass[final]{clv2025}

\jvol{vv}
\jnum{nn}
\jyear{2026}

\dochead{Long Paper}

% The journal fills in the three dates on acceptance. Leaving their template
% placeholder in would print a literal "DD Month YYYY" on page one of the
% preprint, which reads as an unfinished draft.
\pageonefooter{Preprint. Dataset, code and raw model outputs: https://doi.org/10.5281/zenodo.22695091}

\usepackage{amsmath}
\usepackage{booktabs}

% Cyrillic examples appear in Sections 7.3 and 7.4, so T2A is loaded alongside
% the T1 the class already selects. T1 stays the default and \cyr switches
% encoding for the few Cyrillic words, which is why they are all wrapped.
\usepackage[T2A,T1]{fontenc}
\DeclareRobustCommand{\cyr}[1]{{\fontencoding{T2A}\selectfont #1}}

% Azerbaijani `ə' (U+0259) has no T1 slot, and tipa's T3 schwa exists in
% neither the body face nor the typewriter one, so it would substitute a
% mismatched font on the four occurrences of the letter. A rotated `e' in the
% current font is what a schwa is, and it matches whatever face surrounds it.
\DeclareRobustCommand{\azschwa}{\rotatebox[origin=c]{180}{e}}
\DeclareUnicodeCharacter{0259}{\azschwa}

\runningtitle{Script, Not Relatedness}
\runningauthor{Garibli}

% The class prints "(c) <year> Association for Computational Linguistics" at the
% foot of page one. That is true of a published article and false of a preprint
% that has not been accepted, so while this is a preprint the line names the
% author instead. ON ACCEPTANCE, DELETE THE NEXT LINE and the class default
% returns.
% \makeatletter is required: \@jyear carries an @, which is not a letter in the
% document body, so without it the year prints as the literal text "jyear".
\makeatletter
\renewcommand{\footmark}{\copyright\ \@jyear\ Nihat Garibli. Preprint; not peer reviewed.}
\makeatother

\begin{document}

\title{Script, Not Relatedness: What Fine-Tuning\\
on a Neighbouring Language Costs Azerbaijani}

% clv2025.cls builds the page-one footer with \@footnotetext but never sets
% \@thefnmark, so \maketitle fails unless at least one \thanks has run first.
\author{Nihat Garibli\thanks{Corresponding author.}$^{,1}$}

\affilblock{
    \affil{Baku State University\\\quad \email{qaribli.nihat.nemat@bsu.edu.az}}
}

\maketitle

\begin{abstract}
Azerbaijani is not missing from multilingual evaluation: TUMLU contributes 735
items and INCLUDE 6,937. Both are multiple choice, and neither is parallel, so
neither can compare a model's Azerbaijani with its own English on the same
question. The format matters too: of the items a model answers correctly under
multiple choice, only 0.5\% to 10.3\% are also produced correctly in free form,
so selection and production are close to dissociated.

We introduce AZ-Eval, 1006 short-answer items stated in parallel Azerbaijani
and English, of which 708 form a control stratum whose subject matter is
universal by construction, so that failure there is a language failure rather
than a knowledge gap. Human verification is a build gate rather than a label:
the build refuses any row not individually checked. An independent annotator
re-checked a stratified sample of 200 items and agreed on 93.5\%; the 11 items
they rejected were removed, one of which exposed a defective question template.

Measuring 37 models under a single protocol, we find the gap concentrated in
the models that low-resource work actually deploys: 3.9 to 43.5 points across
29 open models of 8B or smaller, against 3.7 to 14.2 points for five large
commercial models. Scale does not explain it, since a 235B model trails four
smaller ones by roughly ten points; nor does quantisation, since an 8B model
served at full precision behaves like its 4-bit peers.

Our main experiment asks what adapting a model to a neighbouring language costs
a third language. Thirteen base/fine-tuned pairs, ten of which pass our
measurement gates, cross two factors: whether the adaptation target is related
to Azerbaijani, and whether it shares its script. Reporting Azerbaijani loss
minus English loss, so that ordinary forgetting is subtracted, relatedness
carries no predictive power and script carries most of it. The closest
relative, Turkish, does no damage in either of two pairs; a wholly unrelated
Slavic language, Russian, damages Azerbaijani in both of its pairs. Every
Latin-script target leaves Azerbaijani intact (four pairs); four of five
Cyrillic-script targets damage it by 8 to 11 points. A non-Latin target is not
by itself sufficient: a model adapted to Thai, Burmese and Tamil alongside four
Latin-script languages, sharing its base with the most damaged pair, costs
Azerbaijani nothing. The tokeniser is measured rather than assumed and is
identical in seven of the ten pairs, so the contrast survives with segmentation
held exactly constant. An exact permutation test over the ten pairs puts the
script account ahead of three rivals, and also shows what the design cannot do:
every damaged pair is built on a Qwen model, so the script of the target and
the family of the base predict these outcomes almost equally well. We name the
one experiment that would separate them.

The single exception, a Ukrainian-adapted model that does no damage, is
absorbed by a marker that separates all ten pairs without exception: whether
the adapted model begins writing Azerbaijani answers in Cyrillic. Every pair
that emits any Cyrillic is damaged; every pair that emits none is not. We stop
short of calling this the mechanism, because transliteration recovers only part
of the Kazakh damage and almost none of the Russian.

The failure decomposes into two modes with opposite prospects. A Kazakh-adapted
model writes Azerbaijani answers in Cyrillic, and transliteration more than
doubles its score: the knowledge is present and the surface form is wrong. A
Turkish-adapted model writes them in Turkish, which no normalisation repairs
because the alphabet is already correct, and which the standard word-level
language-confusion detector misses in 92\% of cases, since the Turkish alphabet
is a subset of the Azerbaijani one and no character ever leaves the expected
range.
\end{abstract}

%%%%%%%%%%%%%%%%%%%%%%%%%%%%%%%%%%%%%%%%%%%%%%%%%%%%%%%%%%%%%%%%%%%%%%%%%%%%%
\section{Introduction}
%%%%%%%%%%%%%%%%%%%%%%%%%%%%%%%%%%%%%%%%%%%%%%%%%%%%%%%%%%%%%%%%%%%%%%%%%%%%%

Azerbaijani is spoken by roughly ten million people and is not absent from
multilingual evaluation: TUMLU contributes 735 items and INCLUDE 6,937. Both,
however, are multiple choice and neither is parallel, so neither can answer the
question this paper asks: how a single model's Azerbaijani compares with its
own English on the same item. Selecting among four options is also not the same
skill as producing an answer: on identical items, only 0.5\% to 10.3\% of what
a model gets right under multiple choice is also produced correctly in free
form.

Its closest well-resourced Turkic relative, Kazakh, has both a benchmark and
several openly released adapted models. That asymmetry suggests an appealing
shortcut: adapt a Kazakh model and inherit its Turkic competence.

We built a benchmark to test that shortcut, and it does not work. Adapting a
model to Kazakh makes it \emph{worse} at Azerbaijani than the model it was
adapted from, beyond what general forgetting explains.

The cause is the script, and relatedness has nothing to do with it. Turkish is
a closer relative than Kazakh and shares the Latin alphabet; adapting to it
costs Azerbaijani nothing. Russian is not related to Azerbaijani at all and is
written in Cyrillic; adapting to it costs 8 to 10 points. Across ten measured
pairs relatedness splits exactly evenly and predicts nothing, while every
Latin-script target is harmless and four of five Cyrillic-script targets are
not.

Establishing that required a benchmark, and building the benchmark produced
findings of its own. This paper reports both.

\subsection{Where the problem actually lives}

We measure 37 models, and the picture is not uniform. Among 29 open models of
8B parameters or smaller, Azerbaijani trails English by 3.9 to 43.5 points,
every gap significant at $p \le 0.0001$. Among five large models accessed
through a commercial API, the gap narrows to between 3.7 and 14.2 points. A
sixth API model, \texttt{qwen3-8b}, is the same size as the local class and
behaves like it (27.5 points) despite full-precision serving.

This is a more precise claim than ``Azerbaijani is poorly supported'', and we
state it precisely because the imprecise version is now false. Three things
follow.

\paragraph{The problem is concentrated in the models people actually deploy
locally} Low-resource language work is done disproportionately with open,
small, quantised models: they run on one GPU, they cost nothing per call, and
they can be used where data cannot leave the building. That is exactly the
class where the gap is 20 to 40 points.

\paragraph{At the very top the gap becomes hard to establish, and we report
that rather than overstate it} For \texttt{gpt-4o} the Azerbaijani-English
difference is 3.7 points with a bootstrap interval of $[1.0, 6.4]$, but it does
not survive Holm correction within its family ($p = 0.10$), and under any
looser normalisation the interval covers zero. We therefore claim a gap for the
\emph{class} of models we measure, not for the single strongest one.

Nor is the human baseline a ceiling. On the 50 items a speaker answered,
\texttt{claude-sonnet-5} scores 74.0\% against the human's 70.0\%. The frontier
is past the point where ``models cannot do Azerbaijani'' is a defensible
summary, which is precisely why the interesting question is where the failure
remains.

\paragraph{Model size does not predict Azerbaijani ability}
\texttt{qwen3-235b-a22b} is larger than every one of the four leading models
yet its gap is three to four times wider. Four models from four independent
organisations converge within 3.1 points of each other, while a much larger
model sits far outside that band. Whatever produces Azerbaijani competence, it
is not parameter count.

\subsection{Contributions}

\begin{enumerate}
\item \textbf{AZ-Eval}, 1006 short-answer items stated in parallel Azerbaijani
and English, to our knowledge the first open parallel, free-form set for the
language, with every item human-verified as a gate rather than a label. An
independent second annotator re-checked a stratified sample of 200 items;
agreement was 93.5\% and the 11 items they rejected were removed.

\item \textbf{A control stratum of 708 items} whose subject matter is universal
by construction. Models demonstrably know these facts in English, so failure in
Azerbaijani is a language failure, not a knowledge gap. The gap is \emph{wider}
on this stratum, not narrower.

\item \textbf{A nine-pair design crossing relatedness with script}, over six
target groups, eight laboratories and seven distinct base models. Relatedness
turns out to carry no predictive power; script carries nine of ten. An exact
permutation test over the ten pairs separates the script account from chance
and ranks it above three competing accounts, including the one the project
started from. The tokeniser is identical within six of the pairs, so the
contrast holds with segmentation exactly constant. Two further pairs were
declared and then excluded by measurement gates, and are reported rather than
dropped. The same test shows that the base model family predicts the outcome
almost as well as the script does, and we report that rather than leave it for
a reader to find.
\namecite{ibrahimzade2026turkic} predicted a script effect for this language
family six months earlier, in a framework paper with no experiments; this is
its test.

\item \textbf{A marker that separates every pair without exception}: whether
the adapted model begins writing Azerbaijani in Cyrillic. All four damaged
pairs do; none of the five undamaged pairs does, including the one
Cyrillic-script target that causes no damage and a model adapted to three
non-Latin scripts that also causes none. We report it as a marker rather than a
mechanism, because transliteration recovers only part of the damage where the
account predicts it should recover all of it. We further show that the standard
word-level language-confusion detector is blind to this pair by construction,
since the Turkish alphabet is a subset of the Azerbaijani one, and measure that
it misses 92\% of these errors.

\item \textbf{The scale result above}, including the negative finding that size
does not predict the gap.

\item \textbf{A set of documented negative results and measurement failures},
because the ones that cost us the most time will cost others the same. One of
them is a measurement that confirmed our hypothesis and was wrong, caught by
reading raw model output rather than by any statistic.
\end{enumerate}

\subsection{What this paper is not}

It is not a claim that frontier models cannot handle Azerbaijani. They handle
it reasonably well, and we report the numbers that show it.

It is not a scaling study. Six API models cannot establish a scaling law, and
the largest of them contradicts the monotone reading anyway.

It is not a general theory of script effects. Our ten pairs rest on eight base
models drawn from only four families, and every pair that shows damage is built
on a Qwen model, so the script of the target and the identity of the base cannot
be told apart here. Section~\ref{sec:between} measures how close the two
accounts are instead of leaving the reader to notice, and the limitations name
the single experiment that would decide between them.

%%%%%%%%%%%%%%%%%%%%%%%%%%%%%%%%%%%%%%%%%%%%%%%%%%%%%%%%%%%%%%%%%%%%%%%%%%%%%
\section{Related Work}
%%%%%%%%%%%%%%%%%%%%%%%%%%%%%%%%%%%%%%%%%%%%%%%%%%%%%%%%%%%%%%%%%%%%%%%%%%%%%

\subsection{Azerbaijani is covered by existing benchmarks, but not in a way
that can measure a language gap}

We must be precise here, because the loose version of our motivation is false.
Azerbaijani is \emph{not} absent from multilingual evaluation.

TUMLU \cite{isbarov2025tumlu} is a natively developed benchmark for eight
Turkic languages, middle- and high-school level, across eleven school subjects,
totalling 38,139 four-choice items. Its Azerbaijani portion contributes 735
items to the manually verified TUMLU-mini subset, which we import and audit.

Our audit is not a criticism of that dataset; their own paper anticipates it,
reporting that in languages such as Azerbaijani, where questions were
community-developed, ``around 10\% of the questions were either invalid or had
incorrect answers''. We retain 521 of 735 after filtering ten defect classes,
which is consistent with that estimate and is why the imported rows carry
\texttt{verified\_by: external} rather than \texttt{human}.

INCLUDE \cite{romanou2025include} covers 44 written languages with 197,243 QA
pairs drawn from local examinations. Azerbaijani is among them, with 6,937
items, listed as Latin script, Turkic, and mid-resource.

KazMMLU \cite{togmanov2025kazmmlu} is not about Azerbaijani but is central to
our motivation: it provides 10,969 Kazakh and 12,031 Russian items. Kazakh has
a resource that Azerbaijani lacked, and open Kazakh-adapted models exist. That
asymmetry is precisely what makes ``adapt a Kazakh model'' an appealing
shortcut, and what this paper tests.

Two properties are shared by all three, and both matter for us.

\emph{They are not parallel.} Items are developed natively or drawn from local
exams, so there is no English counterpart of the same item. A model's
Azerbaijani score can be compared to another model's Azerbaijani score, but not
to its own English score on the same question. The quantity this paper
measures, the within-model within-item language gap, cannot be computed from
them.

\emph{They are multiple choice.} This measures selection among given options,
not production of an answer. The distinction is not cosmetic: on identical
items from TUMLU-az, four models score 34.2 to 38.8 points higher under
multiple choice than under short answer.

That raw difference should not be read on its own, because guessing earns 25\%
under four-way choice and close to nothing under free-form answering. The
chance-immune version of the statistic is retention: of the items a model
answers correctly in multiple choice, only 0.5\% to 10.3\% are also produced
correctly in free form. Recognition and production are close to dissociated
here, and a multiple-choice benchmark therefore overstates usable ability in
the generation setting where these models are actually deployed.

Work on the model side of the same problem is also under way:
\namecite{isbarov2024azerbaijani} release open foundation models adapted to
Azerbaijani, which is the class of model this paper measures rather than
builds. Their effort and ours meet at the same obstacle from opposite
directions, and the pairs in Section 6 are precisely about what such an
adaptation costs a language it was not aimed at.

AZ-Eval is complementary rather than competing. It is smaller than either (1006
items), and it is parallel and free-form, which is what the questions in this
paper require.

\subsection{Script and transliteration in cross-lingual transfer}

A line of work establishes that script divergence obstructs transfer between
related languages, and that transliteration to a shared script recovers part of
it. \namecite{moosa2023transliteration} ask directly whether transliteration
helps multilingual language modelling; \namecite{xhelili2024breaking} propose
transliteration-based post-pretraining alignment and report improvements up to
50\% on some tasks; RomanSetu \cite{husain2024romansetu} romanises inputs to
unlock multilingual capability; and further work analyses why transliteration
improves cross-lingual alignment.

Our question is the mirror image of theirs, and it is the gap we fill. All of
these study the \emph{pre-training} or \emph{alignment} stage and ask whether
converting to a shared script improves transfer \emph{to the target language}.
We ask what fine-tuning on a different-script relative does to a \emph{third}
language, and we measure degradation rather than gain. To our knowledge that
question has not been asked in this literature, which we verified by reading
the alignment paper in full: it reports only improvements and does not measure
harm to any language.

Our transliteration result supports their mechanism from the other side.
\texttt{Qolda-AVL-5B} produces 85.9\% of its Azerbaijani answers in Cyrillic,
and transliteration more than doubles its score: the knowledge is present and
the surface form is wrong.

TUMLU already measures a script effect, on a different axis, and we must not
present ours as the first. They build dual-script versions of their Kazakh,
Crimean Tatar and Uzbek subsets and compare model accuracy across them (their
Table 6). Their finding is that a model reads a language best in whichever
script dominates its training corpus: Crimean Tatar is better in Latin, Kazakh
better in Cyrillic, Uyghur better in Arabic.

Their axis is the script of the \emph{input}, under multiple choice, with the
model unchanged. Ours is the script of the \emph{output}, after adaptation,
under free-form generation. The two are complementary, and theirs supplies the
quantity that makes our mechanism concrete: in FineWeb 2 they count
1,837,049,585 Cyrillic and \emph{zero} Latin Kazakh words. A model adapted to
Kazakh is, in corpus terms, adapted to Cyrillic. That is the sharpest external
support our script explanation has, and it comes from a benchmark paper rather
than from us.

They also observe, as we do, that generation is the harder setting: their
conclusion notes that models understand these languages reasonably well but
``are less capable of generating text in these languages''.

\subsection{The strongest competing explanation: tokenisation, not script}

\namecite{tufa2024unknown} conclude that the tokeniser is a stronger factor
than shared script, language similarity, or model size. Read carelessly, this
contradicts our central claim, and we treat it as the objection to beat rather
than a citation to bury.

Their setup differs from ours in three ways that matter. They use encoder
models (BERT, RoBERTa, ALBERT, Arabic-BERT, mBERT, CANINE) on NER and POS
tagging with Amharic as the target; they measure performance \emph{on the
target language only}, never degradation of a third language after adaptation;
and their tokeniser result emerges from comparisons \emph{across} six models
whose tokenisers differ by construction.

They and we do not measure the same quantity, and saying so is necessary rather
than defensive. Their tokeniser variable is the \emph{algorithm class}: their
Table 1 varies WordPiece, BPE, SentencePiece and character-level tokenisation
across models. Ours is the \emph{vocabulary difference within a single
base-tuned pair}, where the algorithm is identical by construction because the
tuned model inherits its base's tokeniser. Their finding is that the algorithm
class matters across models; ours is that within a pair where the vocabulary is
also unchanged, the script contrast survives. These are compatible claims, and
neither is evidence against the other.

Our design holds the tokeniser constant \emph{within} a pair, and we verify
that it is in fact constant rather than assuming it (Section 6.5). In the two
pairs where it is, the Cyrillic-Latin contrast remains: $+9.8$ points of excess
damage against $+1.5$.

Their finding is therefore the reason our design controls what it controls, not
evidence against the result.

\subsection{Catastrophic forgetting, and how our measure relates to existing
ones}

That fine-tuning degrades previously acquired capabilities is well established
\cite{mccloskey1989catastrophic,kemker2018measuring,luo2023empirical}. The
closest work to ours measures it in the multilingual setting directly:
\namecite{liu2025conditions} fine-tune translation models and plot
per-language-pair performance before against after, and
\namecite{koloski2023measuring} measure forgetting across chains of
cross-lingual transfer.

Normalised forgetting measures already exist, and we do not claim otherwise.
Koloski et al. adopt the \namecite{kemker2018measuring} family, of which the
relevant member is
%
\begin{align}
\Omega_{\mathrm{base}} = \frac{1}{T}\sum_{i=1}^{T}
    \frac{\alpha_{\mathrm{base},i}}{\alpha_{\mathrm{ideal}}}
\end{align}
%
where $\alpha_{\mathrm{ideal}}$ is the model's own offline score. This is a
\emph{ratio of a language to its own ideal}. Our excess damage is a
\emph{difference between two languages} measured on the same items:
%
\begin{align}
\Delta_{\mathrm{excess}} =
    (\mathrm{AZ}_{\mathrm{before}} - \mathrm{AZ}_{\mathrm{after}}) -
    (\mathrm{EN}_{\mathrm{before}} - \mathrm{EN}_{\mathrm{after}}).
\end{align}
%
The distinction is not cosmetic. In Kemker's $\Omega$ and in Koloski et al.,
the reference language is the \emph{source} language, which the model was
trained on. In Liu and Niehues, forgetting is an absolute change per language
pair against that pair's own pre-fine-tuning score. In our design \emph{neither}
language is a fine-tuning target: Azerbaijani and English are both bystanders,
and one bystander calibrates the other. We have not found that construction in
this literature, and we state the claim at that resolution rather than claiming
to have invented forgetting normalisation.

The literature also justifies the choice. Absolute Azerbaijani loss is not
interpretable alone, because some of it is ordinary forgetting; using the
model's own English degradation as the baseline subtracts it. This is what
explains one of our pairs. \texttt{Trendyol} loses 26.9 points in English,
which is textbook catastrophic forgetting rather than Azerbaijani-specific
harm, and is why its excess damage is strongly negative without Azerbaijani
being spared.

\subsection{Language confusion is a known failure with known metrics}

Our second failure mode, an Azerbaijani prompt answered in Turkish, is not a
new observation, and an earlier draft of this paper overstated it by saying no
metric catches it. That is false as stated, and the corrected version is
sharper.

\namecite{marchisio2024language} name the phenomenon \emph{language confusion},
build the Language Confusion Benchmark over 15 languages, and define
line-level and word-level pass rates (LPR, WPR) using off-the-shelf language
identification. Liu and Niehues report LangID accuracy for the same purpose and
observe it collapse to 22.1\% on unseen pairs. The failure mode is established
and measured.

What we add is that these detectors are structurally blind to the Azerbaijani
and Turkish case. For languages in Latin script, Marchisio et al. define
word-level detection as the presence of a character outside the target script's
Unicode range. The Turkish alphabet's 29 letters are a subset of Azerbaijani's
32; Azerbaijani adds only \texttt{ə}, \texttt{x} and \texttt{q}. A word written
in Turkish orthography therefore cannot leave the Azerbaijani range, and the
detector is silent by construction rather than by accident.

We measured this rather than arguing it. Applying the rule to our 159
human-labelled errors, it fires on 2 of the 24 wrong-language rows and misses
92\%. Both rows it catches contain \texttt{w}, a letter in neither alphabet, so
the two hits are English intrusions and not the Turkish confusion at issue.

\begin{table}[h!]
\caption{The word-level language-confusion rule applied to our human-labelled
errors. It is silent by construction on the Azerbaijani-Turkish pair.}
\label{tab:confusion-related}
\centering
\begin{tabular}{lr}
\toprule
 & $n$ \\
\midrule
Human-labelled errors & 159 \\
Labelled wrong-language & 24 \\
Word-level rule fires & 2 \\
Missed & 22 (92\%) \\
\bottomrule
\end{tabular}
\end{table}

Azerbaijani is not among the benchmark's 15 languages, which is consistent with
this being an untested configuration rather than a known limitation.

We do not extend the claim to the line-level rule. That rule splits a response
into lines and checks each with fastText, whereas our answers average 1.34
words; fastText is not installed in our environment and we therefore assert
nothing about how it would behave on this pair.

\subsection{The closest prior work is Turkic-specific, and one paper predicts
our result without testing it}

Three 2026 papers work on cross-lingual transfer within the Turkic family, and
they are the nearest neighbours of this study.

\namecite{ibrahimzade2026turkic} develop a theoretical framework for adaptation
across Azerbaijani, Kazakh, Uzbek, Turkmen and Gagauz. They propose a Turkic
Transfer Coefficient combining morphological similarity, lexical overlap,
syntactic structure and \emph{script compatibility}, and they predict that
``Kazakh demonstrates somewhat lower TTC values with the Latin-script languages
due to the continued use of Cyrillic orthography'', and that ``orthographic
divergence, particularly the coexistence of Latin and Cyrillic scripts, may
partially reduce effective transfer despite high morphological similarity''.

That is close to our hypothesis, in our language family, stated six months
earlier. It is also explicitly untested: they write that ``the goal is not to
report experimental outcomes, but rather to establish a conceptual model'', and
the paper contains no experiments. We therefore present this work as the
empirical test of a published prediction rather than as an independent idea.

The test returns a verdict favourable to the script half of their coefficient
and unfavourable to the typological half. Their TTC combines script
compatibility with morphological and lexical similarity, so it predicts that a
typologically close, script-compatible target transfers best and that
typological distance costs something. Across ten pairs we find that script
compatibility does the work and typological similarity does none: Turkish, the
most similar target available, costs Azerbaijani nothing, while Russian, about
as distant as a target can be, costs 8 to 10 points.

The concrete amendment our measurement supports is therefore to weight the
script term heavily and the typological terms near zero, at least for damage to
a third language. One of our ten pairs (Ukrainian) is script-incompatible and
costs nothing, so even the script term is not deterministic.

\namecite{cinar2026turkic} build pairwise transfer matrices for machine
translation among Turkish, Azerbaijani, Uzbek, Kazakh and Kyrgyz with mT5, and
report that latinisation improves BLEU and chrF when the transfer target is
Cyrillic, by as much as 63.6\% BLEU for Azerbaijani to Kyrgyz. Their
transliteration acts on the \emph{input} of a translation system; ours acts on
the \emph{output} of a generative model that has already been adapted. Their
finding that transfer is strongest for Turkish-Azerbaijani also underwrites our
design choice: Turkish is the closer relative, which is why a relatedness
account predicts the Turkish pairs should be the damaging ones.

\namecite{zhang2026script} compares two script-unification schemes,
general-purpose romanisation against the family-specific Common Turkic Script,
using fastText on eleven Turkic languages evaluated on NER and POS. Like the
transliteration literature above, the aim is to \emph{improve} transfer by
removing script differences, not to measure what script differences destroy.

None of the three measures degradation of a language that is not a target of
the adaptation, and none evaluates free-form generation.

\subsection{A finding that points the other way, and why we report it}

\namecite{khelli2025knowledge} study knowledge loss across 52 languages with
XLM-R and LoRA adapters on MASSIVE slot filling, and conclude that languages
written in Latin script are strong donors and stable receivers, suffering less
catastrophic forgetting than non-Latin ones.

Read naively, this predicts that Azerbaijani, a Latin-script language, should
be relatively safe, and in three or four of our ten pairs it is not. The resolution is
that we are not varying the same thing. Their variable is the script of the
language being \emph{forgotten}; ours is the script of the language being
adapted \emph{to}. A Latin-script bystander enjoys no protection when the
adaptation target is Cyrillic, because the damage arrives through what the
model is being taught to write rather than through what it is being asked to
retain. Their design contains no case where the target and the bystander differ
in script, so the two results do not actually conflict.

Their models are also encoders performing sequence labelling, and Azerbaijani
is not among their 52 languages. We report the tension rather than omit it,
because a reader who knows that paper will otherwise think we do not.

%%%%%%%%%%%%%%%%%%%%%%%%%%%%%%%%%%%%%%%%%%%%%%%%%%%%%%%%%%%%%%%%%%%%%%%%%%%%%
\section{AZ-Eval}
%%%%%%%%%%%%%%%%%%%%%%%%%%%%%%%%%%%%%%%%%%%%%%%%%%%%%%%%%%%%%%%%%%%%%%%%%%%%%

\subsection{Design}

AZ-Eval contains 1006 short-answer items, each stated in parallel Azerbaijani
and English. Parallelism is the property that makes the central measurement
possible: without it there is no within-item comparison, and a language gap
cannot be separated from a difficulty difference.

\begin{table}[h!]
\caption{AZ-Eval at a glance.}
\label{tab:dataset}
\centering
\begin{tabular}{lr}
\toprule
Items & 1006 \\
Control stratum & 708 (70\%) \\
Majority baseline & 0.8\% \\
Distinct answers & 748 under scoring normalisation (769 surface forms) \\
Mean answer length & 1.34 words \\
\bottomrule
\end{tabular}
\end{table}

\paragraph{Categories} Mathematics 264, geography 188, science 179, culture
109, world 104, language 85, history 77.

\paragraph{Provenance} \texttt{wikidata-template} 524,
\texttt{computed-template} 236, \texttt{manual} 246.

Answers are short by design. A one- or two-word answer can be scored by exact
match under an explicit normalisation policy; a sentence cannot, and any
attempt to score sentences would measure instruction-following rather than
knowledge. We return to this in Section 4.1.

\subsection{Human verification is a build gate}

The field \texttt{verified\_by} is not metadata. \texttt{build\_dataset}
refuses to emit any row whose value is not \texttt{human}, so the claim ``every
item was checked'' is enforced by the code that produces the artefact rather
than asserted in the paper.

Rejected drafts are retained under \texttt{data/raw/} with the reason recorded,
because the acceptance rate is itself a quality statistic and discarding
rejections would destroy it.

A separate state, \texttt{external}, marks rows verified by the authors of an
outside dataset and audited by us but not individually re-verified. Such rows
never enter the final dataset. Calling them \texttt{pending} would have hidden
the external verification; calling them \texttt{human} would have claimed work
we did not do.

\subsection{The control stratum}

The most predictable objection to a benchmark for a low-resource language is
that its items are locally specific, so a model fails in that language because
it does not know the fact in any language. The control stratum answers this:
708 items whose subject matter is universal by construction: world capitals and
currencies, chemical elements, and computed mathematics.

Membership is decided by \emph{template}, not by category label. A category
label is too coarse: \texttt{geography} contains both ``the capital of Russia''
and ``the mouth of the Oxchuchay'', and only the first is universal. The
\texttt{country\_capital} template, by contrast, harvests world countries under
a \texttt{sitelinks >= 8} filter and is universal by construction regardless of
the label its rows carry.

One universal template is deliberately excluded.
\texttt{language\_writing\_system} asks which script a language uses; since the
paper's central claim concerns script, including such items in the control
stratum would create an unnecessary coupling. Sixteen items are lost and an
argument is gained.

\subsection{Contamination-proof items}

236 items are produced by algorithm from 20 question families (polygon angles,
$n$th prime, remainder, binary to decimal, and so on), in four syntactic
variants each, and were created \emph{after} every model reported here. They
cannot appear in any training corpus, because they did not exist.

Correctness is verified two independent ways: a second implementation inside
the test suite that deliberately avoids the generator's own method (subtraction
instead of \texttt{\%}, an explicit loop instead of a closed form, bit-by-bit
accumulation instead of \texttt{int(x, 2)}), and separately with \texttt{sympy}.
Both report zero errors.

Answers never appear inside their own question; the generator drops any item
where the answer string occurs in the question text.

Numeric answers bypass the morphological alias generator, which treats every
digit sequence as a year and would accept ``180-ci il'' (``the year 180'') for
a question about degrees.

\subsection{An imported subset, audited}

The Azerbaijani portion of TUMLU-mini (521 items, CC BY 4.0) is imported for
one purpose only: a format contrast between multiple choice and short answer on
identical items. Importing another dataset means importing its defects, so the
rows were audited for ten defect classes (\texttt{comma\_split},
\texttt{broken\_decimal}, \texttt{duplicate\_choice}, \texttt{meta\_answer},
\texttt{answer\_in\_question}, and others). The correct option is placed at
each position in turn, so the uniform answer distribution is a property of
construction rather than luck.

These rows carry \texttt{verified\_by: external} and are not part of the 1006.

\subsection{Quality measured from outside}

A stratified sample of 200 items was re-checked by a second annotator who had
not seen the dataset. Their file carried no trace of the first annotator's
decision: \texttt{verified\_by}, \texttt{notes} and \texttt{source} were
removed, because the template name and source URL would have biased the
judgement. This is enforced by a test.

\begin{table}[h!]
\caption{Independent re-check of a stratified sample of 200 items.}
\label{tab:agreement}
\centering
\begin{tabular}{lr}
\toprule
Agreement & 93.5\% \\
Flagged & 13 (6.5\%) \\
Removed after review & 11 (5.5\%) \\
Retained after review & 2 \\
\bottomrule
\end{tabular}
\end{table}

The number is credible in both directions: zero rejections would indicate an
easy sample or an inattentive pass, and a high rate would indicate a poor
dataset.

One rejection led to a \emph{systematic} fix. The annotator rejected an item
reading ``Find the third angle: 47 degrees and 63 degrees'' on the grounds that
the figure is never named, so the answer is undetermined. Inspection showed
this was not a single bad row but one variant of a generator family, affecting
three items in both languages. The template was corrected.

The statistic is asymmetric and we say so. Every row in the dataset was
accepted by the first annotator, so their label is always ``correct'' and
Cohen's kappa reduces to a function of the second annotator's rejection rate. A
symmetric measurement would require both annotators to judge the same rows from
scratch, blind to each other; that is future work.

%%%%%%%%%%%%%%%%%%%%%%%%%%%%%%%%%%%%%%%%%%%%%%%%%%%%%%%%%%%%%%%%%%%%%%%%%%%%%
\section{Method}
%%%%%%%%%%%%%%%%%%%%%%%%%%%%%%%%%%%%%%%%%%%%%%%%%%%%%%%%%%%%%%%%%%%%%%%%%%%%%

\subsection{Prompting and decoding}

All models receive the same prompt in the language of the question, with two
few-shot examples and an explicit instruction to answer briefly. Decoding is
greedy with a fixed seed and a 32-token budget.

The few-shot examples are not decoration. With instruction alone, one model
answered eight of eight questions correctly in substance while scoring 1/8,
because it produced ``The capital of France is Paris'' rather than ``Paris''.
Such a number measures output format, not knowledge. The two example facts do
not appear in the dataset.

The Azerbaijani and English prompts ask about the same two example facts,
because differing examples would leak a difficulty difference into the language
comparison.

\subsection{Four nested normalisations}

Exact match requires an explicit policy about what counts as the same answer.
Ours is a chain, and each link adds exactly one transformation:

\begin{displaytext}
STRICT $\rightarrow$ MORPH \quad + suffix stripping\\
MORPH $\rightarrow$ LENIENT \quad + diacritic folding\\
LENIENT $\rightarrow$ TRANSLIT \quad + Cyrillic-to-Latin transliteration
\end{displaytext}

Every transformation is applied symmetrically to prediction and gold. Because
the chain is nested, the question ``which step recovered this error'' has
exactly one answer, and that answer identifies the cause. The chain is
therefore the paper's measuring instrument, not merely a cleanup step, and
Section 7 reads it as such.

Accepted-answer sets are additionally expanded by rule-based Azerbaijani
morphological generation, so that inflected forms of a correct answer are
credited.

\subsection{Statistics}

All model comparisons are paired at the item level. We report percentile
bootstrap confidence intervals (1000 resamples, fixed seed) and $p$ values from
a sign-flipping permutation test. Multiple comparisons are corrected by the
Holm procedure \cite{holm1979simple}.

\subsection{Confirmatory and exploratory families are separated}

This is the paper's main methodological commitment and it is stated openly.

The analysis code compares every model pair it encounters and applies Holm
correction across that list. When fourteen models were added for breadth, the
family grew to 474 tests and the $p$ value of the project's founding hypothesis
drifted from 0.018 to 0.047, even though those fourteen models do not test that
hypothesis at all.

Penalising a pre-declared hypothesis for exploratory tests added later is
wrong. The declared pairs therefore form their own family of 104 tests (13
pairs $\times$ 4 chains $\times$ 2 languages), corrected within it. The exploratory
sweep is not deleted: it remains in the tables and is corrected within its own
family.

This is not narrowing a family until a result appears. The pair list is fixed
in the module and in the commit history, and it predates the sweep. When a
fourth pair was added the family grew from 24 to 32 tests, and when a fifth was
added it grew to 40. Each enlargement makes Holm \emph{stricter}: growing a
family can only cost significance, never grant it. The fifth pair was declared
before its result was computed, for a reason given in Section 6.1.

\subsection{Coverage gate}

A run covering less than 90\% of the dataset is excluded from comparison
tables. The reason is that a truncated run is not a random sample: the
evaluation walks the dataset in order, so a 22-row run consists of the first 22
rows, which are template-built capital-city questions and systematically
easier. Such a run scores 83\% where the complete run scores 16\%.

Printing an $N$ column is not sufficient protection, because readers compare
percentages that sit side by side.

\subsection{Degenerate runs are not measurements}

A run in which more than 20\% of raw responses are empty is excluded from the
robustness comparison. If a model emitted nothing, there is nothing to score.

This rule is not tuned to a result. It was introduced after a hybrid reasoning
model, asked without few-shot examples, returned empty content on 42.7\% of
English items, which collapsed its English score and made the language gap
appear to vanish. The gap did not vanish because Azerbaijani improved; it
vanished because the English side broke. Reporting that as evidence about the
hypothesis would have been wrong.

Measured across all runs before the rule was adopted: 53 of 56 had exactly zero
empty responses, one model unable to do the task at all had 79-84\%, and the
degenerate condition had 42.7\%. The gate creates no borderline cases.

The rate is computed on \emph{raw} text rather than extracted answers: an empty
extraction may indicate a defective extractor, whereas empty raw text means the
model produced nothing.

\subsection{Three more ways a run stops being a measurement}

A run in which more than 20\% of responses begin by repeating a few-shot
example answer is likewise excluded. The pathology is that the model emits the
example answers on their own lines and only then the real one, so an extractor
that takes the first line scores a correct answer as wrong.

This gate was added after it cost us a result. One model did this on 67.8\% of
Azerbaijani items and 0\% of English items, which broke one side of a two-sided
comparison and produced 7.0 points of excess damage where the correct value is
0.1. Section~
\ref{sec:failures} reports the episode.

The rule is derived from the production prompt templates rather than restated:
a response counts as an echo when it has more than one line and its first line
appears in the prompt prefix that the run itself was given. Restating the
example answers in the checking code would let a template change slip past the
gate unnoticed.

Measured across all runs, one exceeds the threshold at 68.2\% and the next
highest is 8.7\%, so the gate creates no borderline cases. As with the
empty-response rule, the threshold was chosen from a gap in the data rather
than tuned to a result.

A second variant of the same pathology has to be checked separately. Some
models return the \emph{question} instead of an answer. That is not caught by
the rule above, because such a response is a single line and matches the
question rather than the examples, so we detect it by comparing each response
against the question the item actually posed.

This rule found a run we had already excluded for a different reason. The
\texttt{zeroshot} condition that produced 42.7\% empty English responses
(Section 8.1) also repeats the Azerbaijani question in 31.4\% of items. Two
independent gates rejecting the same run is a useful sign that neither is
arbitrary.

A third gate generalises both. Whatever the model does wrong, the consequence
is the same: the answer is present in the text and does not reach the score. We
therefore measure the share of scored-wrong items whose gold answer appears in
the raw response. Over 106 runs the distribution is one run at 42.4\% and then
a 26-point gap down to 16.4\%, so the threshold sits at 30\%, inside the only
gap the data contains. A first attempt at 15\% was wrong: it fell in the middle
of the dense cluster and made five runs borderline.

Response \emph{length} is deliberately not the rule. \texttt{Mistral-7B-v0.1}
answers in full sentences on 100\% of English items and still scores 52.6\%,
because the extractor handles them. The pathology is a lost answer, not a long
one.

Figure~\ref{fig:gates} shows what each of the five gates actually removed.

\begin{figure}[t!]
\centering
\includegraphics[width=0.72\textwidth]{figures/fig_gates.pdf}
\caption{Runs surviving each measurement gate, applied in order to all 106
runs. Every gate excluded something, which is why none of them is decoration.
The gates are applied before any result is read, not after.}
\label{fig:gates}
\end{figure}

\subsection{Generation is separate from scoring}

\texttt{run\_eval.py} writes raw model output only; all scoring happens in
\texttt{analyze.py}. A change to a scoring rule therefore never requires
re-running a model, and every historical run can be re-scored under a new rule.

\subsection{Two protocol constraints for API models}

Hybrid reasoning models must have reasoning disabled. Otherwise the 32-token
budget is consumed by internal reasoning and the visible answer is empty; we
measured 9 empty responses out of 10 in that condition. The local reasoning
model is run with reasoning disabled for the same reason, so this preserves
comparability rather than saving cost.

Some endpoints do not permit it. Three we intended to include return
\texttt{Reasoning is mandatory for this endpoint and cannot be disabled}.
Raising the token budget for those models alone would have made them
incomparable with the other 28, so they were excluded rather than run under
different conditions. We record this because it is a practical obstacle for
anyone building a short-answer benchmark today, not a property of our setup.

%%%%%%%%%%%%%%%%%%%%%%%%%%%%%%%%%%%%%%%%%%%%%%%%%%%%%%%%%%%%%%%%%%%%%%%%%%%%%
\section{The Azerbaijani-English Gap}
%%%%%%%%%%%%%%%%%%%%%%%%%%%%%%%%%%%%%%%%%%%%%%%%%%%%%%%%%%%%%%%%%%%%%%%%%%%%%

\subsection{The gap is present in every model we measured}

We evaluate 29 open models from 23 independent organisations on all 1006 items,
under identical greedy decoding and 4-bit quantisation. Every model scores
lower in Azerbaijani than in English, and every gap is significant at
$p \le 0.0001$ under a paired sign-flipping permutation test.

\begin{table}[h!]
\caption{Open models, STRICT normalisation, on the full dataset (994 to 1006
items per model, every run above the 90\% coverage gate of Section 4.5). The
English column is printed beside the gap because a small gap in a model that
cannot do the task in English carries no information.
$*$ quoted from its \texttt{oneshot} run, the style in which this model is a
valid measurement (Section~\ref{sec:failures}).
$\ddagger$ fails the empty-response gate at 83.5\% and is kept because its
3.9-point gap is the clearest illustration that a small gap can mean a model
with nothing to lose. Two further open models,
\texttt{Latxa-Qwen3-VL-4B-Instruct} and \texttt{Racka-4B}, were measured and
are omitted: their Azerbaijani runs fail the question-echo gate, so their
scores would report a protocol failure rather than an ability.}
\label{tab:open-models}
\centering
\begin{tabular}{lrrrr}
\toprule
Model & AZ & EN & Gap & 95\% CI \\
\midrule
\texttt{Qwen3.5-4B-Base} & 37.5\% & 60.0\% & 22.4 & [19.5, 25.3] \\
\texttt{Meta-Llama-3-8B} & 33.7\% & 52.6\% & 18.9 & [16.4, 21.5] \\
\texttt{Qwen-SEA-LION-v4-4B-VL} & 28.5\% & 54.8\% & 26.2 & [23.3, 29.2] \\
\texttt{Qwen3-VL-4B-Thinking} & 27.9\% & 55.2\% & 27.4 & [24.2, 30.3] \\
\texttt{Qwen3-VL-4B-Instruct} & 27.1\% & 54.6\% & 27.6 & [24.3, 30.4] \\
\texttt{Turkish-Llama-8b-v0.1} & 23.5\% & 44.0\% & 20.4 & [17.7, 23.3] \\
\texttt{borealis-4b} & 22.6\% & 45.3\% & 22.8 & [19.8, 25.6] \\
\texttt{gemma-3-4b-it} & 20.5\% & 46.9\% & 26.4 & [23.2, 29.8] \\
\texttt{Qwen3-4B} & 20.1\% & 45.4\% & 25.3 & [22.2, 28.4] \\
\texttt{QVikhr-3-4B-Instruction} & 18.3\% & 45.9\% & 27.6 & [24.6, 30.5] \\
\texttt{Mistral-7B-v0.1} & 17.2\% & 52.6\% & 35.4 & [32.5, 38.3] \\
\texttt{RuadaptQwen3-4B-Hybrid} & 17.2\% & 48.5\% & 31.3 & [28.2, 34.2] \\
\texttt{Qwen3.5-4B-Base-Kazakh} & 17.1\% & 50.2\% & 33.1 & [30.0, 36.2] \\
\texttt{Luth-1.7B-Instruct} & 10.7\% & 39.0\% & 28.2 & [25.4, 31.2] \\
\texttt{MamayLM-Gemma-3-4B-IT}$^{*}$ & 20.2\% & 47.0\% & 26.8 & [23.9, 29.5] \\
\texttt{Phi-4-mini-instruct} & 11.2\% & 44.6\% & 33.4 & [30.4, 36.4] \\
\texttt{aya-expanse-8b} & 11.2\% & 47.4\% & 36.2 & [33.1, 39.2] \\
\texttt{Qwen3-1.7B} & 6.5\% & 28.6\% & 22.0 & [19.1, 24.7] \\
\texttt{Trendyol-LLM-7b-base-v1.0} & 6.3\% & 25.8\% & 19.4 & [16.7, 22.2] \\
\texttt{mGPT} & 5.6\% & 14.7\% & 9.1 & [7.0, 11.1] \\
\texttt{Qolda-AVL-5B} & 4.4\% & 41.5\% & 37.1 & [33.8, 40.2] \\
\texttt{Phi-3.5-mini-instruct} & 3.1\% & 46.6\% & 43.5 & [40.3, 46.5] \\
\texttt{next-1b} & 2.7\% & 16.6\% & 13.9 & [11.7, 16.2] \\
\texttt{Falcon3-3B-Instruct} & 1.4\% & 34.3\% & 32.9 & [29.9, 36.0] \\
\texttt{granite-3.1-2b-instruct} & 0.3\% & 32.2\% & 31.9 & [28.9, 34.5] \\
\texttt{stablelm-2-1\_6b-chat} & 0.2\% & 7.6\% & 7.4 & [5.7, 9.1] \\
\texttt{EuroLLM-1.7B-Instruct} & 0.1\% & 17.3\% & 17.2 & [14.9, 19.6] \\
\texttt{SmolLM2-1.7B-Instruct} & 0.0\% & 24.7\% & 24.7 & [22.0, 27.7] \\
\texttt{bloomz-1b7}$^{\ddagger}$ & 0.0\% & 3.9\% & 3.9 & [2.8, 5.1] \\
\bottomrule
\end{tabular}
\end{table}

Two readings of this table are wrong and we rule them out explicitly.

\paragraph{A small gap does not indicate good Azerbaijani support}
\texttt{bloomz-1b7} has the smallest gap in the table (3.9 points) because it
scores 3.9\% in English: it has nothing to lose. The gap is interpretable only
for models that can do the task in English at all, which is why the English
column is printed beside it rather than in a separate table.

\paragraph{The gap is not a small-model artefact} The strongest Azerbaijani
model, \texttt{Qwen3.5-4B-Base}, still loses 22.4 points, and the strongest
English model in the set loses more, not less.

\subsection{The gap is a language gap, not a knowledge gap}

The most natural objection to a benchmark of this kind is that the items are
locally specific, so that a model fails in Azerbaijani because it does not know
the fact in any language. We answer this with a control stratum of 708 items
(70\% of the dataset) whose subject matter is universal by construction: world
facts, chemical elements, and computed mathematics. Models demonstrably know
these facts in English.

On the control stratum the gap does not shrink. It widens.

\begin{table}[h!]
\caption{The gap on all items against the gap on the control stratum, whose
subject matter is universal by construction.}
\label{tab:control}
\centering
\begin{tabular}{lrr}
\toprule
Model & Gap, all items & Gap, control stratum \\
\midrule
\texttt{Phi-3.5-mini-instruct} & 43.5 & \textbf{59.9} \\
\texttt{Qolda-AVL-5B} & 37.1 & \textbf{50.4} \\
\texttt{aya-expanse-8b} & 36.3 & \textbf{48.7} \\
\texttt{Qwen3.5-4B-Base} & 22.4 & \textbf{29.5} \\
\bottomrule
\end{tabular}
\end{table}

The failure is therefore in producing Azerbaijani, not in possessing the
knowledge.

\subsection{The gap does not come from the composition of the dataset}

Mathematics is the largest category (264 of 1006 items, 26\%), and the gap is
widest there, so the headline number could in principle be inflated by
composition. We tested this by removing the largest category and recomputing.

\begin{table}[h!]
\caption{Removing the largest category makes the gap larger, not smaller.}
\label{tab:composition}
\centering
\begin{tabular}{lrr}
\toprule
Subset & $n$ & Mean gap \\
\midrule
all items & 1006 & 24.7 \\
without mathematics & 742 & \textbf{25.6} \\
control stratum & 708 & 33.7 \\
control stratum without mathematics & 464 & \textbf{39.0} \\
\bottomrule
\end{tabular}
\end{table}

Removing mathematics makes the gap larger, not smaller. Mathematics dilutes the
effect rather than driving it, and the reason is mechanical: a mathematics
answer is usually a numeral, and numerals have no alphabet, so the orthographic
failure mode described in Section 7 cannot operate there.

\subsection{The gap narrows sharply with model strength}

All 22 models above are 8B parameters or smaller and were run in 4-bit
quantisation, so the obvious question is whether the gap is an artefact of
small open models. We therefore measured six further models through a
commercial routing API under the identical protocol: same items, same prompt,
same 32-token budget, greedy decoding. Five are substantially larger; the
sixth, \texttt{qwen3-8b}, is deliberately the same size as the local class but
served at full precision, which makes it the control for quantisation.

\begin{table}[h!]
\caption{Models served through a commercial routing API, under the identical
protocol. The best local model is repeated for reference.}
\label{tab:api-models}
\centering
\small
\begin{tabular}{llrrrr}
\toprule
Model & Organisation & AZ & EN & Gap & 95\% CI \\
\midrule
\texttt{gpt-4o} & OpenAI & 66.7\% & 70.4\% & \textbf{3.7} & [1.0, 6.4] \\
\texttt{claude-sonnet-5} & Anthropic & 65.5\% & 71.0\% & \textbf{5.5} & [3.0, 7.8] \\
\texttt{deepseek-v3.2} & DeepSeek & 64.1\% & 70.0\% & \textbf{5.9} & [3.5, 8.3] \\
\texttt{llama-4-maverick} & Meta & 63.6\% & 68.5\% & \textbf{4.9} & [2.6, 7.5] \\
\texttt{qwen3-235b-a22b} & Alibaba & 49.8\% & 64.0\% & 14.2 & [11.1, 17.1] \\
\texttt{qwen3-8b} & Alibaba & 28.1\% & 55.7\% & 27.5 & [24.6, 30.4] \\
\midrule
\texttt{Qwen3.5-4B-Base} (local, best) & Alibaba & 37.5\% & 60.0\% & 22.4 & [19.5, 25.3] \\
\bottomrule
\end{tabular}
\end{table}

Three findings follow, and the second is the one we consider most useful.

\paragraph{The gap narrows sharply, and at the very top we stop being able to
establish it} It falls from 27.5 points at 8B to 3.7 points for
\texttt{gpt-4o}. That last value is where the honest reading matters. Its
bootstrap interval, $[1.0, 6.4]$, excludes zero; its Holm-corrected $p$ within
the exploratory family is 0.10, which does not. Under any looser normalisation
chain the interval covers zero as well ($[-0.6, 4.6]$ at LENIENT).

We therefore do not claim a gap for \texttt{gpt-4o} individually. What the six
models support is a claim about the class: five of the six show Holm-significant
gaps, and the trend across them is clear. Declaring the sixth significant on an
uncorrected interval would contradict the correction policy we set out in
Section 4.4, and we are not willing to apply that policy selectively.

\paragraph{The human baseline is not a ceiling either, and an earlier draft of
this section said it was} On the 50 items answered by a speaker,
\texttt{claude-sonnet-5} scores 74.0\% against the human's 70.0\%. Comparing
\texttt{gpt-4o}'s 66.7\% on 1006 items with a 50-item human figure, as we
previously did, is not a valid comparison in the first place; the matched-item
table is the one to read.

\paragraph{Size alone does not predict the gap} \texttt{qwen3-235b-a22b} is
far larger than any of the four leading models yet its gap is three to four
times wider (14.2 against 3.7-5.9), and it is the only one of the six that
produced empty responses. In the other direction, \texttt{qwen3-8b} at full
precision through the same API shows a 27.5-point gap, close to its 4-bit local
peers, so neither the serving stack nor quantisation accounts for the
difference between the groups. Azerbaijani competence therefore does not arrive
automatically with scale; it appears to depend on what the training mixture
contains. For a practitioner choosing a model for Azerbaijani, parameter count
is a poor proxy.

\paragraph{The narrow gap is not one laboratory's peculiarity} Four models
from four independent organisations, trained on different data, land in a band
of 3.1 points in Azerbaijani (63.6 to 66.7) and 2.2 points on the gap (3.7 to
5.9). Convergence this tight across OpenAI, Anthropic, DeepSeek and Meta is not
plausibly a property of any single training pipeline.

\paragraph{Caveats, stated plainly} These six runs used a routing API, so the
serving stack is not under our control and four of the six are not open
weights; they are reported separately from the open models for that reason.
\texttt{qwen3-235b-a22b} returned empty content on 59 Azerbaijani and 13
English items (5.9\% and 1.3\%); excluding those rows its gap is 10.9 rather
than 14.2, and we report both.

\paragraph{A protocol limitation worth recording} Three endpoints we intended
to include (\texttt{gpt-5}, \texttt{gpt-5-mini},
\texttt{gemini-3.1-pro-preview}) refuse to disable reasoning, returning
\texttt{Reasoning is mandatory for this endpoint and cannot be disabled}. Under
a 32-token budget the entire allowance is consumed by internal reasoning and
the visible answer is empty. Raising the budget for those models alone would
have made them incomparable with the other 28, so they were excluded rather
than run under different conditions. Short-answer benchmarking of such
endpoints is not currently possible without changing the protocol for every
model.

\subsection{The dataset is solvable}

A benchmark that no one can solve does not measure models; it merely holds all
of them near zero. We therefore measured a human baseline on a stratified
sample of 50 items, answered by a speaker who had not seen the dataset, under
the same scoring pipeline as the models and with lookup prohibited.

\begin{table}[h!]
\caption{Human baseline on a stratified sample of 50 items, scored under the
same pipeline as the models. The table is printed in full because showing only
the models below the human would support a claim the data does not.}
\label{tab:human}
\centering
\begin{tabular}{lr}
\toprule
 & Score on the same 50 items (TRANSLIT) \\
\midrule
\texttt{claude-sonnet-5} & \textbf{74.0\%} \\
\textbf{Human} & \textbf{70.0\%} \\
\texttt{deepseek-v3.2} & 68.0\% \\
\texttt{llama-4-maverick} & 66.0\% \\
\texttt{gpt-4o} & 66.0\% \\
\texttt{qwen3-235b-a22b} & 52.0\% \\
\texttt{Qwen3.5-4B-Base} (best local) & 36.0\% \\
six models & 0.0\% \\
\bottomrule
\end{tabular}
\end{table}

The human left 7 of 50 items blank (14\%), and blanks are scored as errors, so
this figure is a lower bound on human performance under the stated rules.

The baseline establishes solvability, not a ceiling, and the table is printed
in full for that reason. One model scores above the human and three more sit
within four points. Showing only the models below the human would have made the
same data support a claim it does not support.

We report a correction here. An earlier version of this baseline was answered
by the dataset's own author, who had seen every item and left 59\% blank,
scoring 26.5\%; from that measurement we had concluded that the dataset was
hard for humans too. That conclusion was an artefact of a compromised
measurement and does not survive an independent annotator.

\paragraph{Limitation} One annotator, 50 items, wide interval. The figure
establishes that the task is solvable; it does not establish a precise ceiling.

\subsection{Multiple choice and free-form answering are close to dissociated}

The imported TUMLU-az rows exist for one measurement: the same 521 items put to
the same models in both formats, under identical decoding. This is what
justifies building a free-form set when multiple-choice sets for the language
already exist.

\begin{table}[h!]
\caption{The same 521 items in both formats. The raw difference overstates the
effect, because four-way choice pays 25\% for guessing.}
\label{tab:format}
\centering
\begin{tabular}{lrrr}
\toprule
Model & MCQ & Short answer & Difference \\
\midrule
\texttt{Qwen3-1.7B} & 37.0\% & 2.9\% & $+34.2$ \\
\texttt{Qwen3-VL-4B-Instruct} & 44.9\% & 7.5\% & $+37.4$ \\
\texttt{Qwen3-VL-4B-Thinking} & 43.2\% & 5.8\% & $+37.4$ \\
\texttt{Qolda-AVL-5B} & 39.0\% & 0.2\% & $+38.8$ \\
\bottomrule
\end{tabular}
\end{table}

All four differences have $p = 0.0001$. The raw difference overstates the
effect and we do not rest the claim on it, because four-way choice pays 25\%
for guessing while free-form answering pays almost nothing.

The chance-immune statistic is retention: among the items a model answers
correctly under multiple choice, how many does it also produce correctly in
free form?

\begin{table}[h!]
\caption{Retention: of the items a model recognises under multiple choice, the
share it also produces correctly in free form.}
\label{tab:retention}
\centering
\begin{tabular}{lrr}
\toprule
Model & Correct under MCQ & Also correct free-form \\
\midrule
\texttt{Qwen3-1.7B} & 193 & 3.6\% \\
\texttt{Qwen3-VL-4B-Instruct} & 234 & 10.3\% \\
\texttt{Qwen3-VL-4B-Thinking} & 225 & 10.2\% \\
\texttt{Qolda-AVL-5B} & 203 & 0.5\% \\
\bottomrule
\end{tabular}
\end{table}

Nine out of ten items a model can recognise, it cannot produce. The two formats
are therefore not two ways of measuring one ability, and a multiple-choice
score is not a usable estimate of what a model will do when asked to answer.

This also bears on \texttt{Qolda-AVL-5B} specifically. Its 0.5\% retention
rises to 3.4\% under transliteration, which is the same orthographic story
reported in Section 7 appearing in an independent dataset.

%%%%%%%%%%%%%%%%%%%%%%%%%%%%%%%%%%%%%%%%%%%%%%%%%%%%%%%%%%%%%%%%%%%%%%%%%%%%%
\section{Relatedness or Script?}
%%%%%%%%%%%%%%%%%%%%%%%%%%%%%%%%%%%%%%%%%%%%%%%%%%%%%%%%%%%%%%%%%%%%%%%%%%%%%

\subsection{The question}

Two open models adapted to Kazakh degrade sharply in Azerbaijani. Kazakh
differs from Azerbaijani in two ways at once: it is a related Turkic language,
and it is written in Cyrillic where Azerbaijani uses Latin. No conclusion about
cause is possible while both factors move together.

We separate them with thirteen declared fine-tuning pairs, ten of which pass our
measurement gates. Each is a base model and a variant adapted from it, and they
cross two factors: whether the adaptation target is genealogically related to
Azerbaijani, and whether it shares its script.

\begin{table}[h!]
\caption{The ten pairs that pass every measurement gate. Turkish is a
\emph{closer} relative than Kazakh while sharing Azerbaijani's script, so a
relatedness account predicts more damage from the Turkish pairs, not less.}
\label{tab:pairs-design}
\centering
\footnotesize
\begin{tabular}{llll}
\toprule
Pair & Base $\rightarrow$ adapted & Target & Rel.\ / script \\
\midrule
Kazakh 1 & \texttt{Qwen3-VL-4B-Thinking} $\rightarrow$ \texttt{Qolda-AVL-5B} & Kazakh & yes / Cyrillic \\
Kazakh 2 & \texttt{Qwen3.5-4B-Base} $\rightarrow$ \texttt{Qwen3.5-4B-Base-Kazakh} & Kazakh & yes / Cyrillic \\
Turkish 1 & \texttt{Mistral-7B-v0.1} $\rightarrow$ \texttt{Trendyol-LLM-7b-base} & Turkish & yes / Latin \\
Turkish 2 & \texttt{Meta-Llama-3-8B} $\rightarrow$ \texttt{Turkish-Llama-8b} & Turkish & yes / Latin \\
Cyrillic 3 & \texttt{gemma-3-4b-it} $\rightarrow$ \texttt{MamayLM-Gemma-3-4B-IT} & Ukrainian & no / Cyrillic \\
Russian 1 & \texttt{Qwen3-4B} $\rightarrow$ \texttt{RuadaptQwen3-4B-Hybrid} & Russian & no / Cyrillic \\
Russian 2 & \texttt{Qwen3-4B} $\rightarrow$ \texttt{QVikhr-3-4B-Instruction} & Russian & no / Cyrillic \\
Latin 3 & \texttt{gemma-3-4b-it} $\rightarrow$ \texttt{borealis-4b} & Norwegian & no / Latin \\
Latin 5 & \texttt{Qwen3-1.7B} $\rightarrow$ \texttt{Luth-1.7B-Instruct} & French & no / Latin \\
SEA & \texttt{Qwen3-VL-4B-Instruct} $\rightarrow$ \texttt{Qwen-SEA-LION-v4-4B-VL} & SE Asian & no / mixed \\
\bottomrule
\end{tabular}
\end{table}

Six target-language groups, eight laboratories, seven distinct base models.

The SEA pair is deliberately mixed rather than clean. SEA-LION is adapted to
Vietnamese, Indonesian, Malay and Filipino (Latin) together with Thai, Burmese
and Tamil (three non-Latin scripts), and it shares its base with Kazakh 1, the
pair with the largest damage in the study. It is the closest thing we have to a
test of whether \emph{any} non-Latin adaptation target is enough.

The last nine pairs were added after the first four were analysed, and we say
so rather than presenting thirteen pre-registered pairs. Each was chosen to
attack a limitation of the previous round: Cyrillic 3 because both Cyrillic
pairs adapted to one language from one laboratory; the two Russian pairs
because Cyrillic 3 turned out to be an unusually gentle adaptation; Latin 3 to
fill the last empty cell; and Latin 4 and Latin 5 to break the confound between
the script of the target and the family of the base
(Section~\ref{sec:between}). Every pair was declared in the module, growing the
confirmatory family from 32 tests to 104, before its result was computed.
Enlarging a Holm family can only cost significance, never grant it, so a pair
added to attack our own design cannot help it by arriving.

\subsection{Measuring excess damage rather than absolute loss}

Any fine-tuning degrades other languages to some extent, so absolute
Azerbaijani loss cannot answer the question. We report \emph{excess damage}:
Azerbaijani loss minus English loss, paired item by item, with a bootstrap
interval. English loss serves as the model's own measure of general forgetting.

\begin{table}[h!]
\caption{Excess damage under the \texttt{default} prompt. Bold marks an
interval that excludes zero \emph{and} an Azerbaijani score that actually fell.
$\dagger$ marks the one pair where those two come apart: Latin 5's interval
excludes zero, but both of its languages improved and English improved more, so
nothing was damaged. Section 6.2 states the condition.}
\label{tab:excess-default}
\centering
\small
\begin{tabular}{lllrrrr}
\toprule
Pair & Related & Script & AZ loss & EN loss & Excess damage & 95\% CI \\
\midrule
Kazakh 1 & yes & Cyrillic & $+23.4$ & $+13.7$ & $\mathbf{+9.8}$ & [$+6.2$, $+13.4$] \\
Kazakh 2 & yes & Cyrillic & $+20.4$ & $+9.8$ & $\mathbf{+10.7}$ & [$+7.0$, $+14.3$] \\
Turkish 1 & yes & Latin & $+10.9$ & $+26.9$ & $-16.0$ & [$-19.2$, $-12.8$] \\
Turkish 2 & yes & Latin & $+10.2$ & $+8.7$ & $+1.5$ & [$-1.0$, $+4.5$] \\
Russian 1 & no & Cyrillic & $+2.9$ & $-3.1$ & $\mathbf{+6.0}$ & [$+3.0$, $+8.9$] \\
Russian 2 & no & Cyrillic & $+1.8$ & $-0.5$ & $+2.3$ & [$-0.6$, $+5.2$] \\
Latin 3 & no & Latin & $-1.1$ & $+2.0$ & $-3.1$ & [$-6.1$, $-0.2$] \\
Latin 5 & no & Latin & $-4.1$ & $-10.3$ & $+6.1^{\dagger}$ & [$+3.2$, $+9.2$] \\
SEA & no & mixed & $-1.1$ & $+0.1$ & $-1.2$ & [$-3.3$, $+0.7$] \\
\midrule
\multicolumn{7}{l}{Cyrillic 3: not measurable under this prompt, see Section~
\ref{sec:failures}} \\
\bottomrule
\end{tabular}
\end{table}

Six of the pairs also have a second measurement under the \texttt{oneshot}
prompt, which is the only one valid for Cyrillic 3.

\begin{table}[h!]
\caption{Excess damage under the \texttt{oneshot} prompt, the only style in
which all six of these pairs are clean.}
\label{tab:excess-oneshot}
\centering
\begin{tabular}{lrrrr}
\toprule
Pair & AZ loss & EN loss & Excess damage & 95\% CI \\
\midrule
Kazakh 1 & $+19.1$ & $+10.0$ & $\mathbf{+9.2}$ & [$+6.0$, $+12.5$] \\
Russian 1 & $+8.6$ & $-1.5$ & $\mathbf{+10.1}$ & [$+7.0$, $+13.0$] \\
Russian 2 & $+7.2$ & $-0.8$ & $\mathbf{+8.0}$ & [$+4.9$, $+10.9$] \\
Cyrillic 3 & $+1.2$ & $+1.1$ & $+0.1$ & [$-2.9$, $+3.2$] \\
Latin 3 & $-0.1$ & $+3.0$ & $-3.1$ & [$-6.0$, $-0.5$] \\
SEA & $-0.7$ & $+0.2$ & $-0.9$ & [$-2.8$, $+1.1$] \\
\bottomrule
\end{tabular}
\end{table}

\paragraph{The metric reads as damage only when Azerbaijani actually fell} A
difference of differences carries no information about which language moved,
and two of our pairs sit at opposite ends of that blind spot.

Turkish 1 is the negative end. Its excess damage is $-16.0$, which invites the
reading that Latin fine-tuning protected Azerbaijani. It did not: Azerbaijani
fell 10.9 points, and the sign is negative only because English fell 26.9.

Latin 5 is the positive end, and it is the reason we state the condition rather
than leave it implicit. Its excess damage is $+6.1$ with an interval well clear
of zero, yet Azerbaijani rose from 6.5\% to 10.7\%. English rose further, from
28.6\% to 38.8\%, and the metric reports the shortfall as though something had
been lost. Nothing was.

Excess damage is therefore interpretable as damage only where the model's
Azerbaijani score fell, and both the tables and the code apply that condition
rather than reading the interval alone. The metric was built for the forgetting
regime, where every fine-tune costs the bystander languages something; a
fine-tune that teaches both of them is outside what it was designed to
describe. Of the ten measured pairs, one falls outside, and it is marked
wherever it appears.

Kazakh 1 anchors both tables and agrees across them ($+9.8$ against $+9.2$), as
does Latin 3 ($-3.1$ in both). The Russian pairs move a great deal between
styles ($+6.0$ to $+10.1$, $+2.3$ to $+8.0$) while keeping their sign, which is
why Section 8.1 states this paper's pair claims at the level of direction and
significance rather than magnitude.

Figure~\ref{fig:forest} shows the same intervals sorted, and
Figure~\ref{fig:losses} shows why a difference metric is needed at all.

\begin{figure}[t!]
\centering
\includegraphics[width=0.78\textwidth]{figures/fig_forest.pdf}
\caption{Excess damage with 95\% bootstrap intervals, under the
\texttt{default} prompt. Every pair whose interval lies above zero has a
Cyrillic-script adaptation target; no Latin-script target does. Eight pairs
appear here rather than nine: Cyrillic 3 is not measurable under this prompt
(Section~
\ref{sec:failures}) and is shown in Figure~\ref{fig:marker} under \texttt{oneshot}.}
\label{fig:forest}
\end{figure}

\begin{figure}[t!]
\centering
\includegraphics[width=0.62\textwidth]{figures/fig_losses.pdf}
\caption{Absolute losses in both languages, under the \texttt{default} prompt.
The dashed diagonal is equal loss; a point above it is Azerbaijani damaged
beyond what the model's own English degradation explains. Turkish 1 sits far
below the diagonal because it loses 26.9 points in English, which is why
absolute Azerbaijani loss cannot answer the question on its own.}
\label{fig:losses}
\end{figure}

\subsection{Relatedness predicts nothing; script predicts nine of ten}

Sorting the same ten pairs by each factor in turn is the whole argument.

\paragraph{By relatedness, the split is exactly even and therefore
uninformative}

\begin{table}[h!]
\caption{The ten pairs sorted by whether the adaptation target is
genealogically related to Azerbaijani.}
\label{tab:by-relatedness}
\centering
\begin{tabular}{lll}
\toprule
 & damages Azerbaijani & does not \\
\midrule
Related target (Turkic) & Kazakh 1, Kazakh 2 & Turkish 1, Turkish 2 \\
Unrelated target & Russian 1, Russian 2 & Cyrillic 3, Latin 3, Latin 5, SEA \\
\bottomrule
\end{tabular}
\end{table}

Two of four each way among related targets, two of five among unrelated ones.
The closest relative we could find, Turkish, does no damage; a completely
unrelated Slavic language, Russian, does. Relatedness carries no predictive
power at all, which is a stronger statement than the one an earlier draft of
this paper made and the opposite of the one before it.

\paragraph{By script, nine of ten fall in line}

\begin{table}[h!]
\caption{The same ten pairs sorted by the script of the adaptation target.}
\label{tab:by-script}
\centering
\small
\footnotesize
\begin{tabular}{lll}
\toprule
 & damages & does not \\
\midrule
Cyrillic target & Kazakh 1, Kazakh 2, Russian 1, Russian 2 & Cyrillic 3 \\
Latin target & & Turkish 1, Turkish 2, Latin 3, Latin 5 \\
Mixed target & & SEA \\
\bottomrule
\end{tabular}
\end{table}

Every Latin-script adaptation leaves Azerbaijani intact. Four of five
Cyrillic-script adaptations damage it, by 8 to 11 points of excess damage
across three target-language and laboratory combinations.
Figure~\ref{fig:design} puts the two factors in one frame.

\begin{figure}[t!]
\centering
\includegraphics[width=0.78\textwidth]{figures/fig_design.pdf}
\caption{The design as measured, under the \texttt{default} prompt, which is
the only style in which every cell of the $2\times2$ is occupied. Damage
follows the column (script of the adaptation target) and not the row
(relatedness). Cyrillic 3 is absent because it is not measurable under this
prompt (Section~
\ref{sec:failures}), and the two Russian values are the \texttt{default} ones
($+6.0$, $+2.3$) rather than the larger \texttt{oneshot} ones reported in
Table~\ref{tab:excess-oneshot}; both styles agree in sign. The figure is
generated from the run data rather than drawn by hand, because which cell
damages Azerbaijani is a result, and that result changed three times as pairs
were added.}
\label{fig:design}
\end{figure}

\paragraph{A non-Latin target is not sufficient, which the SEA pair
establishes} SEA-LION is adapted to Thai, Burmese and Tamil alongside four
Latin-script languages, shares its base with the most damaged pair in the
study, and costs Azerbaijani nothing under either prompt style ($-1.2$,
$[-3.3, +0.7]$ under \texttt{default}; $-0.9$, $[-2.8, +1.1]$ under
\texttt{oneshot}). If merely training on a different-script language were
enough, this pair should show it.

\paragraph{The one exception is real and we do not explain it away} Adapting
to Ukrainian costs Azerbaijani nothing measurable. The most visible difference
is that MamayLM barely moved the model in any direction: its combined absolute
change across both languages is 2.3 points, against 8.4 to 37.8 for every other
pair. But ``the model hardly changed'' is close to restating ``Azerbaijani
hardly changed'', so we record it as an observation rather than an explanation.

\subsection{One marker separates all ten pairs}

There is a variable that does what neither factor above manages: whether the
adapted model starts writing Azerbaijani answers in Cyrillic. We found it by
looking at the damaged pairs, not by predicting it in advance, and we label it
exploratory for that reason: with nine points and a continuous variable, a
split that separates them perfectly is close to guaranteed to exist. What makes
it worth reporting is not the separation but where the exception went, which we
return to below.

\begin{table}[h!]
\caption{Cyrillic emission separates the damaged pairs from the undamaged ones
without exception, including the Ukrainian pair that neither factor explains.
Each pair contributes the excess damage from the prompt style in which it is a
valid measurement, so Kazakh 1 and the two Russian pairs are quoted from
\texttt{oneshot} (Table~\ref{tab:excess-oneshot}) and the rest from
\texttt{default} (Table~\ref{tab:excess-default}). Latin 5 is marked
$\dagger$ because its positive figure is not damage: both of its languages
improved (Section 6.2). It emits no Cyrillic and damages nothing, so the marker
is right about it either way.}
\label{tab:marker}
\centering
\begin{tabular}{lrr}
\toprule
Pair & Cyrillic share of AZ answers & Excess damage \\
\midrule
Kazakh 1 & 84.9\% & $+9.2$ \\
Kazakh 2 & 18.9\% & $+10.7$ \\
Russian 1 & 13.8\% & $+10.1$ \\
Russian 2 & 7.7\% & $+8.0$ \\
\midrule
Latin 5 & 0.5\% & $+6.1^{\dagger}$ \\
SEA & 0.2\% & $-1.2$ \\
Cyrillic 3 & 0.0\% & $+0.1$ \\
Turkish 1 & 0.0\% & $-16.0$ \\
Turkish 2 & 0.0\% & $+1.5$ \\
Latin 3 & 0.0\% & $-3.1$ \\
\bottomrule
\end{tabular}
\end{table}

The separation is complete: every pair that emits Cyrillic above its base rate
is damaged, every pair that does not is not, and the Ukrainian exception stops
being an exception. SEA-LION sits at its base model's own 0.2\%, so it counts
as emitting none. A Cyrillic-script target usually pushes a model into writing
Cyrillic Azerbaijani, but not always, and the damage follows the writing rather
than the target. Figure~\ref{fig:marker} shows the separation under both
usable prompt styles.

\begin{figure}[t!]
\centering
\includegraphics[width=\textwidth]{figures/fig_marker.pdf}
\caption{Excess damage against the share of Azerbaijani answers written in
Cyrillic, under both usable prompt styles. The two panels share one vertical
scale computed over both, so that no outlying pair is silently cropped.}
\label{fig:marker}
\end{figure}

We stop short of calling this the mechanism, because transliteration does not
behave as a purely orthographic account requires. If the damage were only the
wrong alphabet, mapping Cyrillic back to Latin should recover it. It recovers
2.5 of Kazakh 1's 19.1 points, and essentially none of Russian 1's 8.6 (8.6 to
8.4) or Russian 2's 7.6 (unchanged). The Cyrillic answers of the
Russian-adapted models are wrong in content as well as in script.

Cyrillic emission is therefore a reliable \emph{marker} of the damage and not a
sufficient account of it. Section 7 examines what the errors actually contain.

\subsection{Tokenisation is controlled in seven of the ten pairs}

The strongest competing explanation in the literature holds that the tokeniser
matters more than the script for cross-lingual transfer
\cite{tufa2024unknown}. We measured, rather than assumed, whether the tokeniser
is held constant inside each pair.

\begin{table}[h!]
\caption{Tokeniser difference within each pair, measured by set difference on
the tuned model's vocabulary against its base. The distinction is categorical:
control tokens do not change how text is segmented, ordinary word pieces do.}
\label{tab:tokenizer}
\centering
\begin{tabular}{lrrl}
\toprule
Pair & Added word pieces & Added control tokens & Control holds? \\
\midrule
Kazakh 1 & 0 & 3 & \textbf{yes} \\
Kazakh 2 & 16\,000 & 0 & no \\
Turkish 1 & 12\,312 & 0 & no \\
Turkish 2 & 0 & 0 & \textbf{yes} \\
Cyrillic 3 & 0 & 0 & \textbf{yes} \\
Russian 1 & 43\,047 & 21 & no \\
Russian 2 & 0 & 0 & \textbf{yes} \\
Latin 3 & 0 & 0 & \textbf{yes} \\
Latin 5 & 0 & 0 & \textbf{yes} \\
SEA & 0 & 0 & \textbf{yes} \\
\bottomrule
\end{tabular}
\end{table}

The distinction is categorical rather than a threshold. Control tokens such as
\texttt{<|audio\_start|>} do not change how text is segmented; ordinary word
pieces do. Examples of the latter are written here with a visible leading
space, because that space is part of the token:
\texttt{\textvisiblespace Almaty}, \texttt{\textvisiblespace Kazakh},
\texttt{\textvisiblespace KZT}.
\texttt{Qolda-AVL-5B} adds three audio control tokens and no word pieces, so
its text tokenisation is identical to its base.
\texttt{RuadaptQwen3-4B-Hybrid} sits at the other extreme: its vocabulary was
deliberately rebuilt for Russian, which is why it adds 43\,047 pieces.

The argument survives restricted to the seven controlled pairs alone, and this
is the cleanest form of it.

\begin{table}[h!]
\caption{The seven pairs in which the tokeniser is identical within the pair. The
script contrast survives with segmentation held exactly constant.}
\label{tab:controlled}
\centering
\begin{tabular}{lllr}
\toprule
Pair & Script & Tokeniser & Excess damage \\
\midrule
Kazakh 1 & Cyrillic & identical & $\mathbf{+9.2}$ \\
Russian 2 & Cyrillic & identical & $\mathbf{+8.0}$ \\
Cyrillic 3 & Cyrillic & identical & $+0.1$ \\
Turkish 2 & Latin & identical & $+1.5$ \\
Latin 3 & Latin & identical & $-3.1$ \\
Latin 5 & Latin & identical & $+6.1^{\dagger}$ \\
SEA & mixed & identical & $-1.2$ \\
\bottomrule
\end{tabular}
\end{table}

Two Cyrillic targets from different language families (Turkic and Slavic),
different laboratories and different base models both damage Azerbaijani with
the tokeniser held exactly constant, while all three Latin targets and the
mixed target do not. Whatever the tokeniser contributes elsewhere, it is not what
produces this contrast. The two Cyrillic bases are both Qwen models, which we
note rather than overstate as independence.

The three uncontrolled pairs are confirmatory rather than load-bearing. It is
worth noting that two of the three extreme values in the study ($-16.0$ and
$+10.7$) belong to them.

\subsection{Testing the between-pair claim, and the confound it exposes}
\label{sec:between}

Everything above is a tally. The confirmatory family of Section 4.4 corrects 88
tests, and every one of them is \emph{within} a pair: it asks whether a tuned
model differs from its own base. The claim this paper is named for is
\emph{between} pairs, and a count of nine out of ten is not a test of it.

We therefore test it directly. With ten pairs split five against five there are
only 126 ways to assign the script labels, so the permutation distribution is
enumerated exactly rather than sampled. The statistic is the difference in mean
excess damage between the group an account predicts to be damaged and the rest;
we also report a rank version, because Turkish 1 sits at $-16.0$ only because
its English collapsed and on means it carries weight it has not earned.

The same test is applied to the rival accounts, which is the part that matters.

\begin{table}[h!]
\caption{Exact permutation test over the ten pairs, under both rules for
choosing a pair's prompt style. Each row asks whether the pairs one account
predicts to be damaged differ from the rest. The relatedness row is the
hypothesis the project began with.}
\label{tab:between}
\centering
\small
\begin{tabular}{lrrrr}
\toprule
& \multicolumn{2}{c}{\texttt{default} preferred} & \multicolumn{2}{c}{\texttt{oneshot} preferred} \\
\cmidrule(lr){2-3}\cmidrule(lr){4-5}
Account & correct & $p$ & correct & $p$ \\
\midrule
Target is Cyrillic & \textbf{8/10} & 0.071 & \textbf{9/10} & \textbf{0.024} \\
Base model is a Qwen & 7/10 & \textbf{0.024} & 8/10 & \textbf{0.014} \\
Adaptation moved the model a lot & 7/10 & 0.567 & 7/10 & 0.375 \\
Target is related to Azerbaijani & 7/10 & 0.967 & 6/10 & 0.786 \\
\bottomrule
\end{tabular}
\end{table}

Three things follow, and the second is the one we would most want a reader to
carry away.

\paragraph{Relatedness is the worst account of the four} It sits below even a
crude measure of how much the fine-tune moved the model at all, which is the
sharpest form of the paper's negative result. The smallest two-sided $p$ this
design can produce is $1/252 = 0.004$, so the values in
Table~\ref{tab:between} are within one order of the floor: ten pairs cannot
deliver a small $p$ no matter how large the effect.

\paragraph{The base model was confounded with the script, and one pair
partly separates them} Until Latin 5 was added, every pair that damaged
Azerbaijani was built on a Qwen model and no pair built on anything else did,
so ``the target is Cyrillic'' and ``the base is a Qwen'' made the same nine
predictions and scored identically. That is not a result about script; it is a
design that cannot tell two variables apart.

The experiment that separates them is a Qwen base with a Latin-script target,
where the script account predicts no damage and the base account predicts
damage. Latin 5 is that case, and Azerbaijani rose. The script account now
predicts eight of ten under the \texttt{default} rule and nine of ten under
\texttt{oneshot}; the base account, which has to call this pair damaged, gets
seven and eight.

\paragraph{The permutation test still favours the base account, and the reason
is the metric rather than the evidence} On classification the script account
leads. On the test in Table~\ref{tab:between} it does not, because the test
ranks groups by excess damage and Latin 5 carries a large positive value
($+6.1$) while having damaged nothing. Including it in the Qwen group lifts
that group's mean. This is the blind spot of Section 6.2 showing up in a
significance test, and we report both rather than choosing the statistic that
flatters the claim.

One further pair would settle it in either direction: a second Qwen-based Latin
target, or a damaging Cyrillic pair on a non-Qwen base. Of the latter we have
one candidate and it does no damage, which is what the base account predicts.
Two Qwen-based Latin candidates were found and could not be measured
(Section~\ref{sec:failures}).

\paragraph{The tally moves with the prompt-style rule and the sign does not}
Under \texttt{default} the count is seven of nine rather than eight, because
Russian 2 is $+2.3$ there with an interval crossing zero and $+8.0$ under
\texttt{oneshot}. This is why the pair claims in this paper are stated at the
level of direction and significance and not magnitude, and why
Table~\ref{tab:between} reports both rules instead of picking the friendlier
one.

\emph{Source: \texttt{results/tables/between\_pairs.md},
\texttt{src/between\_pairs.py}.}

\subsection{Three measurement failures, and how each was caught}
\label{sec:failures}

The first two share a shape: the measurement broke on \emph{one side} of a
two-sided comparison, which is the only kind of error this design cannot
absorb, because the quantity reported is a difference. The third broke both
sides at once and would have entered the tables as a large negative number.

\paragraph{The chat template asymmetry} Two of the first four pairs were
initially measured with the base model receiving a chat template while the
adapted model, having none, received the raw prompt. A pair comparison is
meaningless unless both halves are asked identically, so we re-ran both base
models without a chat template. The correction moved the two Kazakh pairs in
opposite directions and left the conclusion unchanged. It also contradicted our
expectation: the template was not helping the base models but hurting them,
because neither is instruction-tuned.

\paragraph{The example-echo artefact, which agreed with us} The fifth pair
first produced $+7.0$ points of excess damage, comfortably confirming the
hypothesis it was built to test. It was an artefact.

Under the \texttt{default} prompt, which carries two few-shot examples, MamayLM
reproduces the two example answers on their own lines before giving the real
one, so a question whose answer is \texttt{120} is answered with the two
example answers followed by \texttt{120}. It does this on 67.8\% of Azerbaijani
items and 0\% of English items; the base model does it on neither. Our
extractor takes the first line, so the model was scored wrong on items it had
in fact answered correctly: 55\% of the supposedly lost items carry the gold
answer in the raw text.

The cause is the number of examples, not their content.

\begin{table}[h!]
\caption{Example-echo rate by prompt style for MamayLM. The cause is the number
of examples, not their content.}
\label{tab:echo-style}
\centering
\begin{tabular}{llr}
\toprule
Prompt style & Examples & Example-echo rate, AZ \\
\midrule
\texttt{default} & two & 67.8\% \\
\texttt{plain} & two & 100\% \\
\texttt{oneshot} & one & 0.3\% \\
\texttt{zeroshot} & none & 0\% \\
\bottomrule
\end{tabular}
\end{table}

Both halves of the pair were re-run under \texttt{oneshot}, where all four runs
are clean (no empty responses, at most 0.3\% echo), and the excess damage falls
from $+7.0$ to $+0.1$.

\paragraph{A fourth attempt failed a different gate, and it is worth one
paragraph because it validates the gates rather than the claim} We obtained
access to a Hungarian-adapted model (\texttt{Racka-4B}) built on
\texttt{Qwen3-4B}, the same base as both Russian pairs, which would have given
a same-base Latin against Cyrillic comparison, the one contrast the design
still lacks. Its tokeniser turned out to be interesting on its own: the
vocabulary is the \emph{same size} as the base's but 32\,768 entries differ, so
a size check would have called it unchanged.

Given few-shot examples it continues the pattern by rewriting the prompt rather
than completing it, and it does this under all four prompt styles. The
question-echo rule, written the previous day for an unrelated Georgian model,
catches it at 68.4\% of Azerbaijani items against 1.6\% of English ones.

That rule is coarser than the reason, and the distinction turned out to matter
when a second Latin-target candidate arrived. A model that restates the
question we asked and then answers it has lost nothing: the answer is on the
next line and an extractor can reach it. A model that restates a
\emph{different} question, either a few-shot example or one it invented, has
answered something we did not ask, and no extractor can repair that. Only the
second is fatal, and both models do it one-sidedly: \texttt{Racka-4B} answers a
question that was never asked in 29.5\% of Azerbaijani items against 10.3\% of
English ones, and \texttt{Latxa-Qwen3-VL-4B-Instruct}, a Basque, Galician and
Catalan adaptation of \texttt{Qwen3-VL-4B-Instruct}, in 17.3\% against 0.0\%.
The pair we do admit, Latin 5, sits at 0.0\% in both languages.

That is the same one-sided break as every other failure in this section, and it
is the reason both are excluded. We state it in these terms rather than by
citing a threshold, because three Latin-target candidates have now been
excluded and a reader is entitled to ask whether the gates are filtering in a
direction that suits us.

It also exposed a hole in our own code. The pair report was checking only the
extraction gate, and \texttt{Racka-4B} passes that one (12.0\%). A pair can
fail any gate, so the check now runs all of them; the same fix retroactively
removed the Ukrainian pair's invalid \texttt{default} figure from the table,
which had been sitting there because only one gate was consulted.

\paragraph{We report this at length because the artefact pointed our way} A
wrong number that contradicts the hypothesis gets investigated as a matter of
course; a wrong number that confirms it often does not. The check that caught
this one was not statistical but mechanical: read what the model actually
emitted before trusting what it scored. Had we not, the paper would have
reported a confirmation of a claim its own data refutes.

\paragraph{A third pair was declared and then excluded by a gate, which is what
gates are for} We attempted a Georgian target
(\texttt{Kolkha-Mini-Georgian}, base \texttt{Qwen3-1.7B}) to test whether the
script effect extends past Cyrillic. Its Azerbaijani fell from 6.5\% to 0.1\%
and its English from 28.6\% to 4.4\%, which looks like a model destroyed by
adaptation and would have entered the table as $-17.8$ points of excess damage.

Reading its output shows something else. Asked for the capital of Turkey it
answers ``The capital of Turkey is Ankara.'', which is correct and which our
exact-match scoring rejects. Across its English errors, 42.4\% contain the gold
answer in the raw text; the same figure is 1.4\% to 4.2\% for every other model
in the study, including \texttt{Mistral-7B-v0.1} and \texttt{Trendyol}, which
produce full sentences in 100\% and 99.9\% of items respectively and still
score normally because the extractor handles them.

So response length is not the pathology and we do not gate on it. The gate is
on \emph{recoverable errors}: the share of scored-wrong items whose gold answer
is present in the raw text. Measured over 106 runs, the distribution is 42.4\%
and then a 26-point gap to 16.4\%, so the threshold sits at 30\%, inside the
only gap there is. One run fails it.

The pair remains declared in the module and is reported as excluded rather than
deleted, so the attempt and its reason stay in the record.

%%%%%%%%%%%%%%%%%%%%%%%%%%%%%%%%%%%%%%%%%%%%%%%%%%%%%%%%%%%%%%%%%%%%%%%%%%%%%
\section{What the Failures Look Like}
%%%%%%%%%%%%%%%%%%%%%%%%%%%%%%%%%%%%%%%%%%%%%%%%%%%%%%%%%%%%%%%%%%%%%%%%%%%%%

\subsection{Two failure modes, not one}

``The model is bad at Azerbaijani'' conflates at least two different events,
and separating them is practically consequential because only one of them is
recoverable by post-processing.

\begin{table}[h!]
\caption{The failure modes, and whether normalisation can repair each.}
\label{tab:modes}
\centering
\begin{tabular}{lp{0.44\textwidth}p{0.20\textwidth}}
\toprule
Mode & What happens & Does normalisation help? \\
\midrule
Orthographic & the model knows the fact and writes it in Cyrillic & yes, transliteration recovers it \\
Capability & the model cannot produce the answer at all & no \\
Both at once & the model writes Cyrillic \emph{and} the wrong content & no \\
\bottomrule
\end{tabular}
\end{table}

The third row is the one the Russian pairs forced us to add, and Section 7.3
reports it.

We separate them with four nested normalisation chains, each adding exactly one
transformation and applied symmetrically to prediction and gold:

\begin{displaytext}
STRICT $\rightarrow$ MORPH \quad + suffix stripping\\
MORPH $\rightarrow$ LENIENT \quad + diacritic folding\\
LENIENT $\rightarrow$ TRANSLIT \quad + Cyrillic-to-Latin transliteration
\end{displaytext}

Because the chain is nested, the question ``which step recovered this error''
has an unambiguous answer, and that answer is the cause.

\subsection{The two Kazakh models fail differently}

\begin{table}[h!]
\caption{Azerbaijani score under each link of the normalisation chain. The last
column is what the transliteration step alone recovers. Two models adapted to
the same language by the same laboratory fail in measurably different ways.}
\label{tab:chains}
\centering
\small
\begin{tabular}{lrrrrr}
\toprule
Model & STRICT & MORPH & LENIENT & TRANSLIT & Recovered \\
\midrule
\texttt{Qolda-AVL-5B} & 4.4\% & 4.4\% & 4.4\% & 9.2\% & $\mathbf{+4.7}$ \\
\texttt{Qwen3.5-4B-Base-Kazakh} & 17.1\% & 17.7\% & 17.9\% & 19.2\% & $+1.3$ \\
\bottomrule
\end{tabular}
\end{table}

\texttt{Qolda-AVL-5B} produces 85.9\% of its Azerbaijani answers in Cyrillic;
transliteration more than doubles its score.
\texttt{Qwen3.5-4B-Base-Kazakh} produces 18.9\% in Cyrillic, and
transliteration barely moves it.

\subsection{Cyrillic output is not always an orthographic failure}

The Russian-adapted models complicate the tidy two-mode picture, and the
complication is worth stating because it bears directly on the marker in
Section 6.4.

\begin{table}[h!]
\caption{How much of the Azerbaijani loss transliteration recovers. A purely
orthographic reading predicts recovery roughly in proportion to the Cyrillic
share; that holds partially for the Kazakh model and not at all for the Russian
ones. Losses are STRICT before transliteration and TRANSLIT after it, and the
Cyrillic share is of Azerbaijani answers.}
\label{tab:recovery}
\centering
\footnotesize
\begin{tabular}{lrrrr}
\toprule
Model & Cyrillic share & AZ loss & After translit. & Recovered \\
\midrule
\texttt{Qolda-AVL-5B} (Kazakh) & 84.9\% & $+19.1$ & $+16.6$ & $+2.5$ \\
\texttt{RuadaptQwen3-4B-Hybrid} (Russian) & 13.8\% & $+8.6$ & $+8.4$ & $+0.2$ \\
\texttt{QVikhr-3-4B-Instruction} (Russian) & 7.7\% & $+7.6$ & $+7.6$ & $0.0$ \\
\bottomrule
\end{tabular}
\end{table}

A purely orthographic reading predicts that transliteration recovers roughly
the Cyrillic share. It does so partially for the Kazakh model and not at all
for the Russian ones. Inspecting the Russian errors shows why: the Cyrillic
answers are usually wrong in content too, so restoring the alphabet leaves a
wrong answer in the right alphabet. The model returns \cyr{Москва} where the
gold is \emph{Moskva}, which transliteration does repair, but it also returns
\emph{İstanbul} for a question about Ankara, which nothing repairs.

This is why Section 6.4 calls Cyrillic emission a marker rather than a
mechanism. It separates the damaged pairs from the undamaged ones perfectly,
yet the damage it marks is not, in general, the recoverable kind.
Figure~\ref{fig:recovery} makes the shortfall visible.

\begin{figure}[t!]
\centering
\includegraphics[width=0.72\textwidth]{figures/fig_recovery.pdf}
\caption{Azerbaijani loss before and after transliteration, ordered by Cyrillic
share. A purely orthographic account predicts the two bars should converge as
the Cyrillic share rises. They do not, which is why we report Cyrillic emission
as a marker rather than a mechanism.}
\label{fig:recovery}
\end{figure}

\subsection{An automatic taxonomy, and what it leaves to a human}

The chain also yields an error taxonomy without human judgement: whichever step
repairs an error names its cause. On a stratified sample of 100 errors per
model:

\begin{table}[h!]
\caption{Automatic error taxonomy from the normalisation chain. The residual
column is what no chain step explains and a human must label.}
\label{tab:taxonomy}
\centering
\small
\begin{tabular}{lllr}
\toprule
Model & Target & Orthographic & Left to a human \\
\midrule
\texttt{Qolda-AVL-5B} & Kazakh / Cyrillic & \textbf{95\%} & 4\% \\
\texttt{Qwen3.5-4B-Base-Kazakh} & Kazakh / Cyrillic & 21\% & 61\% \\
\texttt{Turkish-Llama-8b-v0.1} & Turkish / Latin & 0\% & 87\% \\
\texttt{Falcon3-3B-Instruct} & (base) & 1\% & 98\% \\
\bottomrule
\end{tabular}
\end{table}

One bucket deserves comment. Transliterating \cyr{Вашингтон} yields
\emph{vaşington} while the gold answer is \emph{vaşinqton}: Cyrillic
\cyr{г} maps to Latin \emph{g}, but Azerbaijani writes that sound as
\emph{q}. The model knows
the fact and the chain still fails to credit it. Counting such rows as factual
errors would inflate the factual category, so they are reported separately.

\subsection{The human half: a second wrong-language failure}

The machine is silent wherever the chain does not help, which is 87-98\% of
errors for capability-type models. We labelled 159 of those residual errors by
hand, across two models chosen to span the two modes.

\begin{table}[h!]
\caption{Hand labels on 159 residual errors. The Kazakh-adapted model fails by
writing in Cyrillic; the Turkish-adapted model fails by writing in Turkish.}
\label{tab:human-labels}
\centering
\begin{tabular}{lrr}
\toprule
Label & Kazakh-adapted & Turkish-adapted \\
\midrule
factual & 29 (41\%) & 37 (42\%) \\
format & 4 (6\%) & 1 (1\%) \\
\textbf{wrong language} & \textbf{1 (1\%)} & \textbf{23 (26\%)} \\
unrelated & 37 (52\%) & 27 (31\%) \\
\bottomrule
\end{tabular}
\end{table}

This completes the picture. The Kazakh-adapted model fails by writing in
Cyrillic; the Turkish-adapted model fails by writing \emph{in Turkish}
(\emph{yapon dili} answered as \emph{Japonca}, \emph{güləşçi} as \emph{Güreş}).

Both are the same kind of event: the model has the fact and emits it in the
wrong linguistic form. The difference is decisive for measurement.
Transliteration recovers the Cyrillic case; no normalisation can recover the
Turkish case, because the alphabet is already correct. That is a claim about
normalisation, not about detectability, and the next paragraphs separate the
two.

This has a consequence for benchmark design beyond the present study, and it
needs stating carefully, because the strong version is false. The failure mode
itself is known and measurable: \namecite{marchisio2024language} call it
language confusion and define word- and line-level pass rates for it.

What we show is that the standard word-level detector is structurally blind to
this particular pair. It flags, for a Latin-script target, any character
outside that language's Unicode range; but the Turkish alphabet is a subset of
the Azerbaijani one, which adds only \texttt{ə}, \texttt{x} and \texttt{q}. A
Turkish word cannot leave the Azerbaijani range, so the detector is silent by
construction.

Applied to our 159 human-labelled errors, the rule fires on 2 of the 24
wrong-language rows and misses 92\% (Table~\ref{tab:confusion-related}). Both
hits contain \texttt{w}, a letter in neither alphabet, so what it catches is
English intrusion rather than the Turkish confusion.

The practical statement is therefore narrower than ``no metric catches this''
and more useful: exact-match accuracy does not distinguish this failure from a
knowledge failure, transliteration cannot repair it, and the off-the-shelf
script-range detector does not see it. A detector for this pair must compare
against a Turkish lexicon rather than an alphabet.

\paragraph{Disclosure} 24 of these rows were relabelled after the first pass,
from \texttt{format} to \texttt{wrong language}, after the criterion was
tightened: an answer that expresses the same concept in another language is a
language error, not a formatting error, whereas an answer that is simply wrong
remains a factual error even if it contains a Turkish place name. Each
relabelled row retains its original label in a comment field.

%%%%%%%%%%%%%%%%%%%%%%%%%%%%%%%%%%%%%%%%%%%%%%%%%%%%%%%%%%%%%%%%%%%%%%%%%%%%%
\section{Robustness}
%%%%%%%%%%%%%%%%%%%%%%%%%%%%%%%%%%%%%%%%%%%%%%%%%%%%%%%%%%%%%%%%%%%%%%%%%%%%%

Each subsection answers an objection a reader would reasonably raise, and each
is a separate module in the codebase, so the answer is computed rather than
asserted.

\subsection{The results do not depend on the prompt}

Four prompt styles differ structurally rather than lexically, since a
paraphrase would make the check ceremonial.

\begin{table}[h!]
\caption{The four prompt styles. They differ structurally, not lexically.}
\label{tab:styles}
\centering
\begin{tabular}{llr}
\toprule
Style & Examples & Labels \\
\midrule
\texttt{default} & two & yes \\
\texttt{oneshot} & one & yes \\
\texttt{plain} & two & no \\
\texttt{zeroshot} & none & yes \\
\bottomrule
\end{tabular}
\end{table}

Both claims, the language gap and the pair difference, hold in the same
direction with intervals excluding zero across every usable style.

\begin{table}[h!]
\caption{The Kazakh 1 pair difference in Azerbaijani across prompt styles. The
magnitude moves and the sign does not.}
\label{tab:prompts}
\centering
\begin{tabular}{lrr}
\toprule
Style & Pair difference (Kazakh 1, AZ) & 95\% CI \\
\midrule
\texttt{default} & $+23.3$ & [20.5, 26.1] \\
\texttt{plain} & $+19.4$ & [17.0, 22.1] \\
\texttt{oneshot} & $+19.0$ & [16.5, 21.5] \\
\texttt{zeroshot} & \multicolumn{2}{l}{excluded, see below} \\
\bottomrule
\end{tabular}
\end{table}

The magnitude moves and the sign does not. We therefore state the claim at the
level of direction and significance, not magnitude.

The choice of style also moves the between-pair tally, and
Table~\ref{tab:between} reports both rules rather than the friendlier one:
under \texttt{oneshot} the script account is right about eight of the nine
pairs and under \texttt{default} about seven, because Russian 2 crosses zero in
the second. The ordering of the four competing accounts is the same either way.

One condition is excluded, and the reason matters. In \texttt{zeroshot} without
a chat template, the base model returned empty content on 42.7\% of English
items; its English score fell from 55.2\% to 7.1\% and the gap ``disappeared''.
The gap did not disappear because Azerbaijani improved but because the English
side broke. This is a failed measurement, not a failed claim, and conflating
the two would misreport the result. The \texttt{oneshot} style was added
afterwards precisely so that the robustness claim would not rest on two styles
instead of three.

Diagnosing this took a wrong turn worth recording: we first assumed the missing
element was a structural cue (\texttt{Answer:}), which \texttt{zeroshot}
already contains. The common feature of the working styles is the presence of
\emph{examples}.

\subsection{Tokenisation cannot explain the contrast}

The tokeniser measurement is reported in Section 6.5 and is not repeated here.
Its conclusion for robustness is that the whole argument survives restricted to
the six controlled pairs. Two Cyrillic targets damage Azerbaijani with the
tokeniser identical (Kazakh 1 $+9.2$, Russian 2 $+8.0$), from different
language families and different laboratories; two Latin targets do not
(Turkish 2 $+1.5$, Latin 3 $-3.1$); and the one Cyrillic target that does no
damage (Cyrillic 3, $+0.1$) is also controlled, so the exception cannot be
blamed on segmentation either. The three extended-vocabulary pairs are
confirmatory rather than load-bearing, and two of the study's three extreme
values ($-16.0$ and $+10.7$) belong to them.

We previously asserted the general rule that fine-tuning does not change the
tokeniser. Measurement refuted it, and the claim is now computed for every pair
and printed with the result.

\subsection{Benchmark contamination cannot explain the gap}

The 236 algorithmically generated items post-date every model measured here.
The gap persists on that subset for every model able to do the task in English,
with values between 27 and 49 points and $p = 0.0001$.

The claim is precise: these \emph{rows} were not in any training corpus. That
the underlying \emph{facts} appear in training corpora is expected and
necessary, since what we measure is whether that knowledge can be expressed in
Azerbaijani.

Models scoring below 5\% in English are marked, and models between 5\% and
10\% are marked separately as near-threshold. A model with nothing to lose
cannot show a gap, and its near-zero value must not be read as evidence of good
Azerbaijani support. The threshold was fixed before the results were seen and
was not adjusted afterwards.

\subsection{The gap does not come from the composition of the dataset}

Mathematics is the largest category and shows the widest gap, so the headline
number could be inflated by composition. Removing the largest category
\emph{increases} the mean gap from 24.7 to 25.6 points, and on the control
stratum from 33.7 to 39.0 (Table~\ref{tab:composition}).

Mathematics dilutes the effect rather than driving it, for a mechanical reason:
a mathematics answer is usually a numeral, numerals have no alphabet, and the
orthographic failure mode cannot operate there.

\subsection{The dataset is solvable}

A speaker who had not seen the dataset scored 70.0\% on a stratified sample of
50 items, against 36.0\% for the best \emph{local} model on the same items,
under the same scoring pipeline and with lookup prohibited. Blanks (14\%) count
as errors, so the figure is a lower bound.

The baseline establishes solvability, not a ceiling, and the distinction is not
rhetorical: \texttt{claude-sonnet-5} reaches 74.0\% on those same items. A
benchmark that no human can pass would not measure models, which is what this
check is for; it was never intended to bound what a model can achieve.

We correct an earlier claim here. A first baseline, answered by the dataset's
own author, scored 26.5\% with 59\% blanks, and from it we had concluded that
the dataset was hard for humans as well. That conclusion was an artefact of a
compromised measurement: the annotator had written the items and was applying a
stricter blank-if-unsure policy than a naive respondent would.

%%%%%%%%%%%%%%%%%%%%%%%%%%%%%%%%%%%%%%%%%%%%%%%%%%%%%%%%%%%%%%%%%%%%%%%%%%%%%
\section{Negative Results}
%%%%%%%%%%%%%%%%%%%%%%%%%%%%%%%%%%%%%%%%%%%%%%%%%%%%%%%%%%%%%%%%%%%%%%%%%%%%%

We report these because the failures cost us the most time, and because a
reader's next idea is likely to be one of them.

\subsection{Prompting for the Latin alphabet does not work}

Instructed explicitly to answer in the Azerbaijani Latin alphabet,
\texttt{Qolda-AVL-5B} barely complies: the Cyrillic share falls from 91.7\% to
88.8\%, and of the seventeen answers that change, one becomes correct.

The behaviour is therefore in the weights, not at the prompt level. This
matters practically: a deployment cannot be fixed by prompt engineering, and
transliteration as post-processing is the cheaper remedy.

\subsection{Auditing gold answers by model agreement does not work}

If several models agree on an answer that differs from the gold, that gold
answer might be wrong. We implemented this, with a permutation test to
establish the chance ceiling: four Qwen models produced 36 suspect items
against a ceiling of 12.

Every suspect we checked by hand had a \emph{correct} gold answer.

The reason is instructive. A permutation test controls for \emph{chance}
agreement, not for \emph{correlated} agreement. Four models from one family
share training data and share mistakes, so they agree on wrong answers far more
often than independence predicts. The method requires genuinely independent
model families, and the statistic that appeared to validate it was measuring
the wrong null.

\subsection{Searching for more Cyrillic pairs, and what the first search cost}

Our Cyrillic evidence originally rested on two pairs adapting to one language
from one laboratory. Our first search for a third concluded that the obstacle
was structural: nearly every Cyrillic adaptation we examined extended the
tokeniser, which breaks the control of Section 6.5, so we stopped looking.

That conclusion was wrong, and stopping was the most expensive mistake in the
project. A systematic search over fine-tunes derived from bases we had already
measured produced four usable pairs within an hour, three of them with an
identical tokeniser. They doubled the design and twice overturned its central
claim.

The reasoning error is worth naming, because it is general. We had two gaps: an
uncontrolled tokeniser and a single target language. We ranked them by how easy
each was to describe rather than by how much each threatened the conclusion.
The single-language gap was much the larger threat, and a pair with an extended
tokeniser would still have been worth measuring.

Two candidates were rejected on inspection rather than on results, and both
rejections mattered. \texttt{Zubr1.0-VL-4B} states in its own model card that
its weights are unchanged from the base: it is a re-host, not an adaptation,
and would have produced a guaranteed null that looked like evidence.
\texttt{BgGPT-7B} (Bulgarian) is a genuine adaptation on a base we had already
measured, but no 7B model would load in our environment; the original
\texttt{Mistral-7B-v0.1} segfaults there too, so this is an environment limit
rather than a property of the model, and the Bulgarian cell remains unfilled.

\subsection{A Georgian pair we declared and could not use}

We attempted a Georgian target to test whether the script effect reaches past
Cyrillic. \texttt{Kolkha-Mini-Georgian} writes Georgian script in 17.0\% of its
Azerbaijani answers, which is exactly the contamination the marker predicts,
and its scores collapse in both languages.

It is not a measurement. The model answers in full sentences (``The capital of
Turkey is Ankara.''), and 42.4\% of its scored-wrong English items contain the
gold answer in the raw text, against 1.4\% to 4.2\% for every other model here.
The score reflects our extractor, not the model's knowledge.

We report the attempt for two reasons. It is the third distinct way a run has
stopped being a measurement in this project, after empty responses and prompt
echo, and the gate that catches it was written for this case and then applied
to all 106 runs. And the pair would have entered the tables at $-17.8$ points
of excess damage, a number with the wrong sign for our own claim, which is a
useful reminder that gates have to run before the result is read, not after.

%%%%%%%%%%%%%%%%%%%%%%%%%%%%%%%%%%%%%%%%%%%%%%%%%%%%%%%%%%%%%%%%%%%%%%%%%%%%%
\section{Limitations}
%%%%%%%%%%%%%%%%%%%%%%%%%%%%%%%%%%%%%%%%%%%%%%%%%%%%%%%%%%%%%%%%%%%%%%%%%%%%%

\begin{enumerate}
\item \textbf{The base model is only partly separated from the script.} Every
pair that damages Azerbaijani is built on a Qwen model, so for most of this
study ``the target is Cyrillic'' and ``the base is a Qwen'' made the same
predictions and scored identically. Latin 5, a French adaptation of
\texttt{Qwen3-1.7B}, is the one pair that pulls them apart: the script account
expects no damage there and is right, the base account expects damage and is
wrong. That is one pair carrying the separation, and two further candidates for
the same cell could not be measured. The permutation test also still favours
the base account, because it ranks groups by excess damage and Latin 5 scores
high on that metric while having damaged nothing
(Section~\ref{sec:between}). We therefore treat the base family as a live
alternative rather than a closed question, and results in Section 6 are best
read as holding for Qwen-derived models until a damaging Cyrillic pair on
another base exists.

\item \textbf{One Cyrillic-script pair does no damage and we cannot say why.}
Adapting to Ukrainian leaves Azerbaijani intact while adapting to Kazakh or
Russian does not. The visible difference is that MamayLM barely moved the model
at all (2.3 points of combined absolute change against 8.4 to 37.8 elsewhere),
but that restates the result rather than explaining it. Until a second
Ukrainian pair or a second gentle Cyrillic pair exists, ``four of five'' is the
honest form of the claim. A Bulgarian pair (\texttt{BgGPT-7B}) would have
helped and could not be loaded in our environment.

\item \textbf{The script contrast is Cyrillic against everything else.} Damage
is demonstrated for Kazakh and Russian targets, absence of damage for Turkish,
Norwegian, Ukrainian and a mixed South-East Asian target. We have no clean pair
for a third script in either direction. A Georgian attempt was declared and
excluded by a measurement gate (Section 9.4), a Hungarian one on the same base
as both Russian pairs was excluded by another, and Greek candidates exist only
at 7B and larger, which our environment cannot load. The same-base Latin
against Cyrillic comparison therefore remains open. ``Script matters''
therefore means ``Cyrillic targets damage Azerbaijani and Latin ones do not'',
and the SEA pair shows that a non-Latin target is not by itself enough.

\item \textbf{Tokenisation is controlled in seven of ten pairs.} The design
survives because the contrast is visible among the controlled pairs alone, but
the three extended-vocabulary pairs carry less weight than their numbers
suggest.

\item \textbf{Local models are 8B or smaller and 4-bit quantised.} Pair
comparisons are internally controlled, so quantisation cannot explain a
within-pair difference, but absolute scores are not comparable to
full-precision serving.

\item \textbf{API models were served through a router.} The serving stack is
outside our control, four of the six are not open weights, and results may
shift as providers change. They are reported separately for this reason.

\item \textbf{The human baseline is one person on 50 items.} It establishes
solvability, not a ceiling. The interval is wide and the respondent was an
educated adult rather than a representative speaker.

\item \textbf{Inter-annotator agreement is asymmetric.} Every row in the
dataset was accepted by the first annotator, so the statistic measures a
rejection rate rather than two independent judgement distributions.

\item \textbf{The error taxonomy's human half covers two models and 159
items.} It is sufficient to contrast the two failure modes but not to estimate
population proportions.

\item \textbf{Short answers only.} The results do not transfer to open-ended
generation or multi-step reasoning, where formatting and length interact with
scoring differently.

\item \textbf{A \texttt{difficulty} field was removed rather than fixed.} It
looked like a per-item judgement but was a template-level constant: across 26
template groups only one carried more than a single value, and rows without one
were silently defaulted to ``medium'', which is why 881 of 1006 read as medium.
Since the value is recoverable from the template recorded in \texttt{notes}, it
carried no information and could only mislead. Raw files may still contain it;
the built dataset does not.
\end{enumerate}

\section*{Data and Code Availability}

The dataset, the evaluation pipeline, the raw model responses for every run
reported here, and the scripts that regenerate every table and figure are
archived at \texttt{https://doi.org/10.5281/zenodo.22695091} and developed at
\texttt{https://github.com/nihatgaribli/AZ-Eval-benchmark}. The DOI resolves to
the current version; the release backing this paper is v2.0.1.

Generation is separate from scoring (Section 4.8), so the archived raw
responses can be re-scored under a different normalisation, a different
extractor or a different gate without re-running a single model. The dataset
under \texttt{data/} and the results under \texttt{results/} are CC BY 4.0; the
pipeline under \texttt{src/} is MIT.

\begin{acknowledgments}
We thank the second annotator, who re-checked a stratified sample of 200 items
and the 50-item human baseline without having seen the dataset beforehand.
\end{acknowledgments}

\bibliographystyle{compling}
\bibliography{references}

\end{document}
